% toolsoftrade/toolsoftrade.tex
% mainfile: ../perfbook.tex
% SPDX-License-Identifier: CC-BY-SA-3.0

\QuickQuizChapter{chp:Tools of the Trade}{常用工具}{qqztoolsoftrade}
%
\Epigraph{你的水平取决于你的工具，而你的工具的水平也取决于你。}{佚名}

本章简要介绍并行编程的一些基本工具，主要关注在类 Linux 操作系统上运行的
用户态应用程序所能使用的工具。
\Cref{sec:toolsoftrade:Scripting Languages}从脚本语言开始，
\cref{sec:toolsoftrade:POSIX Multiprocessing}
描述 POSIX API 支持的多进程并行机制并简要介绍 POSIX thread，
\cref{sec:toolsoftrade:Alternatives to POSIX Operations}
介绍其他环境中的类似操作，最后，
\cref{sec:toolsoftrade:The Right Tool for the Job: How to Choose?}
帮助选择合适的工具来完成任务。

\QuickQuiz{
	你管这些叫工具？？？
	在我看来它们更像是底层同步原语！
}\QuickQuizAnswer{
	它们看起来像底层同步原语，是因为它们实际上就是底层同步原语。
	而且它们确实是构建底层并发软件的基本工具。
}\QuickQuizEnd

请注意，本章仅提供简要介绍。
更多细节可以从参考文献（以及互联网）中获取，
后续章节也将提供更多信息。

\section{脚本语言}
\label{sec:toolsoftrade:Scripting Languages}
%
\epigraph{至高的卓越在于简洁。}
	 {Henry Wadsworth Longfellow，简化版}
% The original:
% "In character, in manner, in style, in all things, the supreme
% excellence is simplicity."

Linux shell 脚本语言提供了简单而有效的并行管理方式。
例如，假设你有一个程序 \co{compute_it}，
需要使用两组不同的参数运行两次。
可以使用 UNIX shell 脚本如下完成：

\input{CodeSamples/toolsoftrade/parallel@compute_it.fcv}

\begin{figure}
\centering
\resizebox{3in}{!}{\includegraphics{toolsoftrade/shellparallel}}
\caption{并行 Shell 执行的执行图}
\label{fig:toolsoftrade:Execution Diagram for Parallel Shell Execution}
\end{figure}

\begin{fcvref}[ln:toolsoftrade:parallel:compute_it]
\Clnref{comp1,comp2}启动该程序的两个实例，
将它们的输出分别重定向到两个不同的文件，
其中 \co{&} 字符指示 shell 在后台运行这两个程序实例。
\Clnref{wait}等待两个实例都完成，
\clnref{cat1,cat2}显示它们的输出。
\end{fcvref}
由此产生的执行过程如
\cref{fig:toolsoftrade:Execution Diagram for Parallel Shell Execution}
所示：
\co{compute_it} 的两个实例并行执行，
\co{wait} 在两者都完成后结束，然后两个
\co{cat} 实例顺序执行。
% @@@ Maui scheduler, load balancing, BOINC, and so on.

\QuickQuizSeries{%
\QuickQuizB{
	但这个愚蠢的 shell 脚本根本不是一个\emph{真正的}并行程序！
	为什么要为这种琐事费心呢？？？
}\QuickQuizAnswerB{
	因为你\emph{永远}不应该忘记简单的东西！

	请记住，本书的标题是
	"并行编程难吗？如果是，你能怎么办？"。
	你能做的最有效的事情之一就是
	不要忘记简单的东西！
	毕竟，如果你选择了困难的方式来做并行编程，
	你只能怪自己。
}\QuickQuizEndB
%
\QuickQuizE{
	有没有更简单的方法来创建并行 shell 脚本？
	如果有，怎么做？
	如果没有，为什么？
}\QuickQuizAnswerE{
	一个直接的方法是 shell 管道：

\begin{VerbatimU}
grep $pattern1 | sed -e 's/a/b/' | sort
\end{VerbatimU}

	对于足够大的输入文件，
	\co{grep} 将与 \co{sed} 编辑操作以及 \co{sort}
	的输入处理并行地进行模式匹配。
	请参阅文件 \path{parallel.sh}，了解
	shell 脚本并行和管道的演示。
}\QuickQuizEndE
}

再举一个例子，\co{make} 软件构建脚本语言提供了一个 \co{-j} 选项，
用于指定在构建过程中引入多少并行度。
因此，在构建 Linux 内核时输入 \co{make -j4}，
可以指定最多同时执行四个构建步骤。

希望这些简单的例子能让你相信，并行编程不一定总是复杂或困难的。

\QuickQuiz{
	但如果基于脚本的并行编程这么简单，为什么还要费心用其他方法？
}\QuickQuizAnswer{
	事实上，目前使用中的并行程序很可能有很大一部分是基于脚本的。
	然而，基于脚本的并行确实有其局限性：
	\begin{enumerate}
	\item	创建新进程通常相当重量级，
		涉及开销很大的 \co{fork()} 和 \co{exec()} 系统调用。
	\item	数据共享（包括管道）通常涉及
		开销很大的文件 I/O。
	\item	脚本可用的可靠同步原语
		通常也涉及开销很大的文件 I/O。
	\item	脚本语言通常太慢，但在协调用较低级编程语言
		编写的长时间运行程序的执行时，
		往往非常有用。
	\end{enumerate}
	这些限制要求基于脚本的并行使用粗粒度并行，
	每个工作单元的执行时间至少为数十毫秒，
	最好更长。

	需要更细粒度并行的人应该认真思考他们的问题，
	看看是否可以用粗粒度的形式来表达。
	如果不行，他们应该考虑使用其他并行编程环境，
	例如
	\cref{sec:toolsoftrade:POSIX Multiprocessing}中讨论的那些。
}\QuickQuizEnd

\section{POSIX 多处理}
\label{sec:toolsoftrade:POSIX Multiprocessing}
%
\epigraph{骆驼是委员会设计出来的马。}{佚名}

本节粗略介绍 POSIX 环境，包括
pthreads~\cite{OpenGroup1997pthreads}，
因为该环境容易获取且被广泛实现。
\Cref{sec:toolsoftrade:POSIX Process Creation and Destruction}
简要展示 POSIX \co{fork()} 及相关原语，
\cref{sec:toolsoftrade:POSIX Thread Creation and Destruction}
简述 thread 的创建和销毁，
\cref{sec:toolsoftrade:POSIX Locking}简要概述
POSIX locking，最后，
\cref{sec:toolsoftrade:POSIX Reader-Writer Locking}描述一种
特定的 lock，可用于被许多 thread 读取但仅偶尔更新的数据。

\subsection{POSIX 进程的创建和销毁}
\label{sec:toolsoftrade:POSIX Process Creation and Destruction}

进程使用 \apipx{fork()} 原语创建，可以使用 \apipx{kill()} 原语销毁，
也可以使用 \apipx{exit()} 原语自行销毁。
执行 \co{fork()} 原语的进程被称为新创建进程的"父"进程。
父进程可以使用 \apipx{wait()} 原语等待其子进程。

请注意，本节中的示例相当简单。
使用这些原语的实际应用程序可能需要操作信号、文件描述符、
共享内存段以及许多其他资源。
此外，某些应用程序需要在特定子进程终止时采取特定操作，
并且可能还需要关注子进程终止的原因。
这些问题当然会给代码增加相当大的复杂性。
更多信息请参阅有关该主题的众多教科书~\cite{WRichardStevens1992,StewartWeiss2013UNIX}。

\begin{listing}
\begin{fcvlabel}[ln:toolsoftrade:forkjoin:main]
\begin{VerbatimL}[commandchars=\%\[\]]
pid = fork();%lnlbl[fork]
if (pid == 0) {%lnlbl[if]
	/* child */%lnlbl[child]
} else if (pid < 0) {%lnlbl[else]
	/* parent, upon error */%lnlbl[errora]
	perror("fork");
	exit(EXIT_FAILURE);%lnlbl[errorb]
} else {
	/* parent, pid == child ID */%lnlbl[parent]
}
\end{VerbatimL}
\end{fcvlabel}
\caption{使用 \tco{fork()} 原语}
\label{lst:toolsoftrade:Using the fork() Primitive}
\end{listing}

\begin{fcvref}[ln:toolsoftrade:forkjoin:main]
如果 \co{fork()} 成功，它会返回两次，一次在父进程中，
一次在子进程中。
\co{fork()} 的返回值允许调用者区分二者，如
\cref{lst:toolsoftrade:Using the fork() Primitive}
(\path{forkjoin.c})所示。
\Clnref{fork}执行 \co{fork()} 原语，并将其返回值
保存在局部变量 \co{pid} 中。
\Clnref{if}检查 \co{pid} 是否为零，如果是，
则这是子进程，继续执行\clnref{child}。
如前所述，子进程可以通过 \co{exit()} 原语终止。
否则，这是父进程，在\clnref{else}检查 \co{fork()} 原语
是否返回了错误，如果是，则在\clnrefrange{errora}{errorb}
打印错误信息并退出。
否则，\co{fork()} 已成功执行，父进程因此执行\clnref{parent}，
此时变量 \co{pid} 包含子进程的进程 ID。
\end{fcvref}

\begin{listing}
\input{CodeSamples/api-pthreads/api-pthreads@waitall.fcv}
\caption{使用 \tco{wait()} 原语}
\label{lst:toolsoftrade:Using the wait() Primitive}
\end{listing}

\begin{fcvref}[ln:api-pthreads:api-pthreads:waitall]
父进程可以使用 \apipx{wait()} 原语来等待其子进程完成。
然而，使用该原语比其 shell 脚本对应物稍微复杂一些，
因为每次调用 \co{wait()} 只等待一个子进程。
因此，通常将 \co{wait()} 封装到一个类似于
\cref{lst:toolsoftrade:Using the wait() Primitive}
(\path{api-pthreads.h})所示的 \apipx{waitall()} 函数中，
该 \co{waitall()} 函数的语义类似于 shell 脚本的 \co{wait} 命令。
\clnrefrange{loopa}{loopb}的循环每执行一次就等待一个子进程。
\Clnref{wait}调用 \co{wait()} 原语，它会阻塞直到一个子进程退出，
并返回该子进程的进程 ID。
如果进程 ID 为 $-1$，则表示 \co{wait()} 原语无法等待子进程。
如果是这样，\clnref{ECHILD}检查 \co{ECHILD} errno，
它表示没有更多的子进程，因此\clnref{break}退出循环。
否则，\clnref{perror,exit}打印错误信息并退出。
\end{fcvref}

\QuickQuiz{
	为什么这个 \co{wait()} 原语需要这么复杂？
	为什么不让它像 shell 脚本的 \co{wait} 那样工作呢？
}\QuickQuizAnswer{
	一些并行应用程序需要在特定子进程退出时采取特殊操作，
	因此需要逐个等待每个子进程。
	此外，一些并行应用程序需要检测子进程死亡的原因。
	正如我们在\cref{lst:toolsoftrade:Using the wait() Primitive}中
	看到的，从 \co{wait()} 函数构建 \co{waitall()} 函数并不难，
	但反过来是不可能的。
	一旦关于特定子进程的信息丢失了，就再也找不回来了。
}\QuickQuizEnd

\begin{listing}
\input{CodeSamples/toolsoftrade/forkjoinvar@main.fcv}
\caption{通过 \tco{fork()} 创建的进程不共享内存}
\label{lst:toolsoftrade:Processes Created Via fork() Do Not Share Memory}
\end{listing}

\begin{fcvref}[ln:toolsoftrade:forkjoinvar:main]
至关重要的是要注意，父进程和子进程\emph{不}共享内存。
这可以通过
\cref{lst:toolsoftrade:Processes Created Via fork() Do Not Share Memory}
(\path{forkjoinvar.c})所示的程序来说明，
其中子进程在\clnref{setx}将全局变量 \co{x} 设置为 1，
在\clnref{print:c}打印一条消息，并在\clnref{exit:s}退出。
父进程继续执行到\clnref{waitall}，在那里等待子进程，
然后在\clnref{print:p}发现它自己的变量 \co{x} 副本仍然为零。
因此输出如下：
\end{fcvref}

\begin{VerbatimU}
Child process set x=1
Parent process sees x=0
\end{VerbatimU}

\QuickQuiz{
	\co{fork()} 和 \co{wait()} 的内容不是比这里讨论的多得多吗？
}\QuickQuizAnswer{
	确实如此，而且本节很可能在未来的版本中扩展，
	以包括消息传递功能（如 UNIX 管道、TCP/IP 和
	共享文件 I/O）以及内存映射（如 \co{mmap()} 和 \co{shmget()}）。
	在此期间，有许多教科书详细介绍了这些原语，
	真正有动力的人可以阅读 man 手册页、使用这些原语的现有并行应用程序，
	以及 Linux 内核实现本身的源代码。

	需要注意的是，
	\cref{lst:toolsoftrade:Processes Created Via fork() Do Not Share Memory}
	中的父进程会等到子进程终止后才执行 \co{printf()}。
	从多个进程并发地使用 \co{printf()} 的缓冲 I/O 写入同一文件
	是非常复杂的，最好避免这样做。
	如果你确实需要并发缓冲 I/O，请查阅你的操作系统文档。
	对于 UNIX/Linux 系统，Stewart Weiss 的讲义提供了
	很好的介绍和有启发性的
	示例~\cite{StewartWeiss2013UNIX}。
}\QuickQuizEnd

最细粒度的并行需要共享内存，这在
\cref{sec:toolsoftrade:POSIX Thread Creation and Destruction}中介绍。
也就是说，共享内存并行可能比 fork-join 并行复杂得多。

\subsection{POSIX Thread 的创建和销毁}
\label{sec:toolsoftrade:POSIX Thread Creation and Destruction}

\begin{fcvref}[ln:toolsoftrade:pcreate:mythread]
要在现有进程内创建一个 thread，可以调用
\apipx{pthread_create()} 原语，例如，如
\cref{lst:toolsoftrade:Threads Created Via pthread-create() Share Memory}
(\path{pcreate.c})的\clnref{create:a,create:b}所示。
第一个参数是指向 \co{pthread_t} 的指针，用于存储待创建 thread 的 ID，
第二个 \co{NULL} 参数是指向可选 \co{pthread_attr_t} 的指针，
第三个参数是新 thread 将要调用的函数（在本例中为 \co{mythread()}），
最后一个 \co{NULL} 参数是将传递给 \co{mythread()} 的参数。
\end{fcvref}

\begin{listing}
\input{CodeSamples/toolsoftrade/pcreate@mythread.fcv}
\caption{通过 \tco{pthread_create()} 创建的 Thread 共享内存}
\label{lst:toolsoftrade:Threads Created Via pthread-create() Share Memory}
\end{listing}

在本例中，\co{mythread()} 只是简单地返回，但它也可以调用
\apipx{pthread_exit()}。

\QuickQuiz{
	如果
	\cref{lst:toolsoftrade:Threads Created Via pthread-create() Share Memory}
	中的 \co{mythread()} 函数可以简单地返回，为什么还要用
	\co{pthread_exit()}？
}\QuickQuizAnswer{
	在这个简单的例子中，完全没有理由。
	然而，想象一个更复杂的例子，其中 \co{mythread()} 调用了
	其他函数，可能是单独编译的。
	在这种情况下，\co{pthread_exit()} 允许这些其他函数
	结束 thread 的执行，而不必将某种错误返回值
	一路传回到 \co{mythread()}。
}\QuickQuizEnd

\begin{fcvref}[ln:toolsoftrade:pcreate:mythread]
\apipx{pthread_join()} 原语（如\clnref{join}所示）类似于
fork-join 的 \apipx{wait()} 原语。
它会阻塞，直到 \co{tid} 变量指定的 thread 完成执行，
无论是通过调用 \apipx{pthread_exit()} 还是从
thread 的顶层函数返回。
thread 的退出值将通过传给 \co{pthread_join()} 的第二个参数
所指向的位置来存储。
thread 的退出值是传给 \co{pthread_exit()} 的值
或 thread 顶层函数的返回值，取决于该 thread 如何退出。
\end{fcvref}

\cref{lst:toolsoftrade:Threads Created Via pthread-create() Share Memory}
所示程序的输出如下，证明了两个 thread 之间确实共享内存：

\begin{VerbatimU}
Child process set x=1
Parent process sees x=1
\end{VerbatimU}

注意，该程序精心确保在任何时刻只有一个 thread 向变量 \co{x} 存储值。
如果一个 thread 可能在向某个变量存储值的同时，
另一个 thread 正在从该变量加载或向其存储，
这种情况被称为\emph{\IX{data race}}。
由于 C 语言不保证 data race 的结果会有任何合理性，
我们需要某种方式来安全地并发访问和修改数据，
例如下一节讨论的 locking 原语。

你说你的 data race 是良性的？
好吧，也许是。
但请帮所有人（包括你自己）一个大忙，
认真阅读\cref{sec:toolsoftrade:Shared-Variable Shenanigans}。
随着编译器的优化越来越激进，
真正良性的 data race 越来越少了。

\QuickQuiz{
	如果 C 语言在存在 data race 时不做任何保证，
	那为什么 Linux 内核有这么多 data race？
	你是想告诉我 Linux 内核完全是坏的吗？？？
}\QuickQuizAnswer{
	啊，但 Linux 内核是用 C 语言的一个精心选择的超集编写的，
	其中包含特殊的 GNU 扩展（如 asm），
	这些扩展即使在存在 data race 的情况下也允许安全执行。
	此外，Linux 内核不在一些 data race 特别成问题的
	平台上运行。
	例如，考虑具有 32 位指针和 16 位总线的嵌入式系统。
	在这样的系统上，涉及对给定指针的存储和加载的 data race
	很可能导致加载返回指针旧值的低 16 位
	与指针新值的高 16 位的拼接。

	尽管如此，即使在 Linux 内核中，data race 也可能相当危险，
	应在可行的情况下
	避免~\cite{JonCorbet2012ACCESS:ONCE}。
}\QuickQuizEnd

\subsection{POSIX 锁}
\label{sec:toolsoftrade:POSIX Locking}

POSIX 标准允许程序员通过"POSIX locking"来避免 data race。
POSIX locking 具有多种原语，其中最基本的是
\apipx{pthread_mutex_lock()} 和 \apipx{pthread_mutex_unlock()}。
这些原语操作的 lock 类型为 \apipx{pthread_mutex_t}。
这些 lock 可以静态声明并使用
\apipx{PTHREAD_MUTEX_INITIALIZER} 初始化，
也可以动态分配并使用 \apipx{pthread_mutex_init()} 原语初始化。
本节的演示代码将采用前一种方式。

\co{pthread_mutex_lock()} 原语"获取"指定的 lock，
\co{pthread_mutex_unlock()} "释放"指定的 lock。
因为这些是"互斥" locking 原语，
在任何给定时刻，只有一个 thread 可以"持有"给定的 lock。
例如，如果一对 thread 同时尝试获取同一个 lock，
其中一个将首先被"授予" lock，
另一个将等待，直到第一个 thread 释放该 lock。
一个简单且相当有用的编程模型允许仅在持有相应 lock 时才能访问
给定的数据项~\cite{Hoare74}。

\QuickQuiz{
	如果我想让多个 thread 同时持有同一个 lock 怎么办？
}\QuickQuizAnswer{
	你首先应该问自己为什么要这样做。
	如果答案是"因为我有很多数据被许多 thread 读取，
	但只是偶尔更新"，那么 POSIX reader-writer lock
	可能是你要找的。
	这些在
	\cref{sec:toolsoftrade:POSIX Reader-Writer Locking}中介绍。

	获得多个 thread 持有同一个 lock 效果的另一种方法是，
	让一个 thread 获取 lock，
	然后使用 \co{pthread_create()} 创建其他 thread。
	至于为什么这会是一个好主意，留给读者思考。
}\QuickQuizEnd

\begin{listing}
\ebresizeverb{.91}{
\input{CodeSamples/toolsoftrade/lock@reader_writer.fcv}
}
\caption{互斥 Lock 的演示}
\label{lst:toolsoftrade:Demonstration of Exclusive Locks}
\end{listing}

\begin{fcvref}[ln:toolsoftrade:lock:reader_writer]
这种互斥 locking 特性通过
\cref{lst:toolsoftrade:Demonstration of Exclusive Locks}
(\path{lock.c})所示的代码来演示。
\Clnref{lock_a}定义并初始化了一个名为 \co{lock_a} 的 POSIX lock，
\clnref{lock_b}类似地定义并初始化了一个名为 \co{lock_b} 的 lock。
\Clnref{x}定义并初始化了一个共享变量~\co{x}。
\end{fcvref}

\begin{fcvref}[ln:toolsoftrade:lock:reader_writer:reader]
\Clnrefrange{b}{e}定义了一个函数 \co{lock_reader()}，
它在持有 \co{arg} 指定的 lock 时反复读取共享变量 \co{x}。
\Clnref{cast}将 \co{arg} 转换为指向 \co{pthread_mutex_t} 的指针，
这是 \co{pthread_mutex_lock()} 和 \co{pthread_mutex_unlock()}
原语所要求的。
\end{fcvref}

\QuickQuizSeries{%
\QuickQuizB{
	为什么不直接让
	\cref{lst:toolsoftrade:Demonstration of Exclusive Locks}
	\clnrefr{ln:toolsoftrade:lock:reader_writer:reader:b}上
	\co{lock_reader()} 的参数是指向 \co{pthread_mutex_t} 的指针？
}\QuickQuizAnswerB{
	因为我们需要将 \co{lock_reader()} 传递给
	\co{pthread_create()}。
	虽然我们可以在传递给 \co{pthread_create()} 时对函数进行类型转换，
	但函数类型转换比简单的指针类型转换丑陋得多，
	也更难正确完成。
}\QuickQuizEndB
%
\QuickQuizE{
	\begin{fcvref}[ln:toolsoftrade:lock:reader_writer]
	\cref{lst:toolsoftrade:Demonstration of Exclusive Locks}的
	\clnref{reader:read_x,writer:inc}上的 \apik{READ_ONCE()}
	和\clnref{writer:inc}上的 \apik{WRITE_ONCE()} 是什么？
	\end{fcvref}
}\QuickQuizAnswerE{
	这些宏约束编译器，以防止它执行对并发访问的
	共享变量有问题的优化。
	它们完全不约束 CPU，除了防止对给定单一变量的
	访问重排序之外。
	注意，这种单一变量约束确实适用于
	\cref{lst:toolsoftrade:Demonstration of Exclusive Locks}
	中所示的代码，因为只访问了变量 \co{x}。

	有关 \co{READ_ONCE()} 和 \co{WRITE_ONCE()} 的更多信息，
	请参阅
	\cref{sec:toolsoftrade:Atomic Operations (gcc Classic)}。
	有关多个 thread 对多个变量的访问排序的更多信息，
	请参阅
	\cref{chp:Advanced Synchronization: Memory Ordering}。
	同时，\co{READ_ONCE(x)} 与 \GCC\ 内置函数
	\co{__atomic_load_n(&x, __ATOMIC_RELAXED)} 非常相似，
	\co{WRITE_ONCE(x, v)} 与 \GCC\ 内置函数
	\co{__atomic_store_n(&x, v, __ATOMIC_RELAXED)} 非常相似。
}\QuickQuizEndE
}

\begin{fcvref}[ln:toolsoftrade:lock:reader_writer:reader]
\Clnrefrange{acq:b}{acq:e}获取指定的
\co{pthread_mutex_t}，检查错误并在出错时退出程序。
\Clnrefrange{loop:b}{loop:e}反复检查 \co{x} 的值，
每次值改变时打印新值。
\Clnref{sleep}休眠一毫秒，使该演示在单处理器机器上也能正常运行。
\Clnrefrange{rel:b}{rel:e}释放 \co{pthread_mutex_t}，
同样检查错误并在出错时退出程序。
最后，\clnref{return}返回 \co{NULL}，同样是为了匹配
\co{pthread_create()} 所要求的函数类型。
\end{fcvref}

\QuickQuiz{
	每次获取和释放 \co{pthread_mutex_t} 都要写四行代码，
	这看起来很痛苦！
	有没有更好的方法？
}\QuickQuizAnswer{
	确实有！
	因此，\co{pthread_mutex_lock()} 和
	\co{pthread_mutex_unlock()} 原语通常被封装在执行错误检查的函数中。
	后面我们将用 Linux 内核的
	\co{spin_lock()} 和 \co{spin_unlock()} API 来封装它们。
}\QuickQuizEnd

\begin{fcvref}[ln:toolsoftrade:lock:reader_writer:writer]
\cref{lst:toolsoftrade:Demonstration of Exclusive Locks}的
\Clnrefrange{b}{e}展示了 \co{lock_writer()}，
它在持有指定 \co{pthread_mutex_t} 时定期更新共享变量 \co{x}。
与 \co{lock_reader()} 类似，\clnref{cast}将 \co{arg}
转换为指向 \co{pthread_mutex_t} 的指针，
\clnrefrange{acq:b}{acq:e}获取指定的 lock，
\clnrefrange{rel:b}{rel:e}释放它。
持有 lock 期间，\clnrefrange{loop:b}{loop:e}递增共享变量 \co{x}，
每次递增之间休眠五毫秒。
最后，\clnrefrange{rel:b}{rel:e}释放 lock。
\end{fcvref}

\begin{listing}
\input{CodeSamples/toolsoftrade/lock@same_lock.fcv}
\caption{使用相同互斥 Lock 的演示}
\label{lst:toolsoftrade:Demonstration of Same Exclusive Lock}
\end{listing}

\begin{fcvref}[ln:toolsoftrade:lock:same_lock]
\Cref{lst:toolsoftrade:Demonstration of Same Exclusive Lock}
展示了一个代码片段，它将 \co{lock_reader()} 和
\co{lock_writer()} 作为 thread 运行，使用相同的 lock，即 \co{lock_a}。
\Clnrefrange{cr:reader:b}{cr:reader:e}创建一个运行
\co{lock_reader()} 的 thread，然后
\clnrefrange{cr:writer:b}{cr:writer:e}创建一个运行
\co{lock_writer()} 的 thread。
\Clnrefrange{wait:b}{wait:e}等待两个 thread 都完成。
该代码片段的输出如下：
\end{fcvref}

\begin{VerbatimU}
Creating two threads using same lock:
lock_reader(): x = 0
\end{VerbatimU}

因为两个 thread 使用相同的 lock，\co{lock_reader()} thread
无法看到 \co{lock_writer()} 持有 lock 时产生的 \co{x} 的
任何中间值。

\QuickQuiz{
	"x = 0"是
	\cref{lst:toolsoftrade:Demonstration of Same Exclusive Lock}
	所示代码片段唯一可能的输出吗？
	如果是，为什么？
	如果不是，还可能出现什么输出，为什么？
}\QuickQuizAnswer{
	不是。
	输出"x = 0"的原因是 \co{lock_reader()} 先获取了 lock。
	如果 \co{lock_writer()} 先获取了 lock，
	那么输出将是 \mbox{``x = 3''}。
	然而，因为代码片段先启动了 \co{lock_reader()}，
	并且此次运行是在多处理器上执行的，
	通常预期 \co{lock_reader()} 会先获取 lock。
	然而，并没有保证，特别是在繁忙的系统上。
}\QuickQuizEnd

\begin{listing}
\input{CodeSamples/toolsoftrade/lock@diff_lock.fcv}
\caption{使用不同互斥 Lock 的演示}
\label{lst:toolsoftrade:Demonstration of Different Exclusive Locks}
\end{listing}

\Cref{lst:toolsoftrade:Demonstration of Different Exclusive Locks}
展示了一个类似的代码片段，但这次使用不同的 lock：
\co{lock_reader()} 使用 \co{lock_a}，
\co{lock_writer()} 使用 \co{lock_b}。
该代码片段的输出如下：

\begin{VerbatimU}
Creating two threads w/different locks:
lock_reader(): x = 0
lock_reader(): x = 1
lock_reader(): x = 2
lock_reader(): x = 3
\end{VerbatimU}

因为两个 thread 使用不同的 lock，它们不会互相排斥，
可以并发运行。
因此 \co{lock_reader()} 函数可以看到
\co{lock_writer()} 存储的 \co{x} 的中间值。

\QuickQuizSeries{%
\QuickQuizB{
	使用不同的 lock 可能会导致相当大的混乱，
	thread 会看到彼此的中间状态。
	那么良好编写的并行程序是否应该限制自己只使用
	一个 lock 来避免这种混乱？
}\QuickQuizAnswerB{
	虽然有时可以使用单个全局 lock 编写既高性能又可扩展的程序，
	但这样的程序是例外而非常规。
	通常需要使用多个 lock 来获得良好的性能和可扩展性。

	此规则的一个可能例外是"事务内存"，
	这目前是一个研究课题。
	事务内存的语义可以粗略地理解为单个全局 lock 的语义，
	同时允许优化并增加了回滚功能~\cite{HansJBoehm2009HOTPAR}。
}\QuickQuizEndB
%
\QuickQuizM{
	在
	\cref{lst:toolsoftrade:Demonstration of Different Exclusive Locks}
	所示的代码中，\co{lock_reader()} 是否保证能看到
	\co{lock_writer()} 产生的所有值？
	为什么？
}\QuickQuizAnswerM{
	不能。
	在繁忙的系统上，\co{lock_reader()} 可能在
	\co{lock_writer()} 的整个执行期间被抢占，
	在这种情况下，它将\emph{看不到} \co{lock_writer()}
	对 \co{x} 的任何中间状态。
}\QuickQuizEndM
%
\QuickQuizE{
	等一下！！！
	\Cref{lst:toolsoftrade:Demonstration of Same Exclusive Lock}
	没有初始化共享变量~\co{x}，
	那为什么在
	\cref{lst:toolsoftrade:Demonstration of Different Exclusive Locks}
	中需要初始化它？
}\QuickQuizAnswerE{
	请参阅\cref{lst:toolsoftrade:Demonstration of Exclusive Locks}的
	\clnrefr{ln:toolsoftrade:lock:reader_writer:x}。
	因为
	\cref{lst:toolsoftrade:Demonstration of Same Exclusive Lock}
	中的代码先运行，它可以依赖于~\co{x} 的编译时初始化。
	\cref{lst:toolsoftrade:Demonstration of Different Exclusive Locks}
	中的代码后运行，所以它必须重新初始化~\co{x}。
}\QuickQuizEndE
}

虽然 POSIX 互斥 locking 还有很多内容，但这些原语提供了
一个良好的起点，而且在许多情况下实际上是足够的。
下一节简要介绍 POSIX reader-writer locking。

\subsection{POSIX 读写锁}
\label{sec:toolsoftrade:POSIX Reader-Writer Locking}

POSIX API 提供了一种 reader-writer lock，
由 \apipx{pthread_rwlock_t} 表示。
与 \apipx{pthread_mutex_t} 一样，\co{pthread_rwlock_t} 可以通过
\apipx{PTHREAD_RWLOCK_INITIALIZER} 静态初始化，
也可以通过 \apipx{pthread_rwlock_init()} 原语动态初始化。
\apipx{pthread_rwlock_rdlock()} 原语以读模式获取
指定的 \apipx{pthread_rwlock_t}，
\apipx{pthread_rwlock_wrlock()} 原语以写模式获取它，
\apipx{pthread_rwlock_unlock()} 原语释放它。
在任何给定时刻，只有一个 thread 可以以写模式持有给定的
\apipx{pthread_rwlock_t}，
但多个 thread 可以以读模式持有给定的 \apipx{pthread_rwlock_t}，
至少在当前没有 thread 以写模式持有它时是这样。

正如你可能预期的，\IXhpl{reader-writer}{lock}是为读多写少的场景设计的。
在这些场景中，reader-writer lock 可以提供比\IXh{exclusive}{lock}
更好的可扩展性，因为 exclusive lock 根据定义在任何给定时刻只允许
一个 thread 持有 lock，
而 reader-writer lock 允许任意数量的 reader 并发持有 lock。
然而，在实践中，我们需要知道 reader-writer lock 能提供多少
额外的可扩展性。

\begin{listing}
\input{CodeSamples/toolsoftrade/rwlockscale@reader.fcv}
\caption{测量 Reader-Writer Lock 的可扩展性}
\label{lst:toolsoftrade:Measuring Reader-Writer Lock Scalability}
\end{listing}

\begin{fcvref}[ln:toolsoftrade:rwlockscale:reader]
\Cref{lst:toolsoftrade:Measuring Reader-Writer Lock Scalability}
(\path{rwlockscale.c})
展示了一种测量 reader-writer lock 可扩展性的方法。
\Clnref{rwlock}展示了 reader-writer lock 的定义和初始化，
\clnref{holdtm}展示了控制每个 thread 持有 reader-writer lock
时间的 \co{holdtime} 参数，
\clnref{thinktm}展示了控制释放 reader-writer lock 到下一次获取之间
时间的 \co{thinktime} 参数，
\clnref{rdcnts}定义了 \co{readcounts} 数组，
每个 reader thread 将其获取 lock 的次数存入其中，
\clnref{nrdrun}定义了 \co{nreadersrunning} 变量，
用于确定所有 reader thread 是否已开始运行。

\Clnrefrange{goflag:b}{goflag:e}定义了 \co{goflag}，
用于同步测试的开始和结束。
该变量初始设置为 \co{GOFLAG_INIT}，在所有 reader thread 启动后
设置为 \co{GOFLAG_RUN}，最后设置为 \co{GOFLAG_STOP} 以终止测试运行。
\end{fcvref}

\begin{fcvref}[ln:toolsoftrade:rwlockscale:reader:reader]
\Clnrefrange{b}{e}定义了 \co{reader()}，即 reader thread。
\Clnref{atmc_inc}原子地递增 \co{nreadersrunning} 变量，
表示该 thread 现在正在运行，
\clnrefrange{wait:b}{wait:e}等待测试开始。
\apik{READ_ONCE()} 原语强制编译器在循环的每次迭代中都获取 \co{goflag}
——否则编译器有权假设 \co{goflag} 的值永远不会改变。
\end{fcvref}

\QuickQuizSeries{%
\QuickQuizB{
	与其到处使用 \apik{READ_ONCE()}，为什么不直接在
	\cref{lst:toolsoftrade:Measuring Reader-Writer Lock Scalability}
	的\clnrefr{ln:toolsoftrade:rwlockscale:reader:goflag:e}
	上将 \co{goflag} 声明为 \co{volatile}？
}\QuickQuizAnswerB{
	在这个特定情况下，\co{volatile} 声明确实是一个合理的替代方案。
	然而，使用 \apik{READ_ONCE()} 的好处是向读者清楚地标明
	\co{goflag} 受到并发读写的影响。
	注意，\apik{READ_ONCE()} 在大多数访问受 lock 保护
	（因此\emph{不会}发生变化）但少数访问在 lock 之外进行的
	情况下特别有用。
	在这种情况下使用 \co{volatile} 声明会使读者更难注意到
	lock 之外的特殊访问，
	也会使编译器更难在 lock 下生成好的代码。
}\QuickQuizEndB
%
\QuickQuizM{
	\apik{READ_ONCE()} 只影响编译器，不影响 CPU。
	难道我们不需要 memory barrier 来确保
	\cref{lst:toolsoftrade:Measuring Reader-Writer Lock Scalability}
	中 \co{goflag} 值的变化及时传播到 CPU 吗？
}\QuickQuizAnswerM{
	不需要，memory barrier 在这里既不需要也没有帮助。
	Memory barrier 只在多个内存引用之间强制排序：
	它们绝对不保证加速数据从系统的一个部分传播到另一个部分。\footnote{
		一直有关于某些硬件中 memory barrier 确实能加速数据传播的
		传闻，但没有确认的目击报告。}
	这引出一条快速经验法则：
	除非你使用多个变量在多个 thread 之间通信，
	否则不需要 memory barrier。

	但是 \co{nreadersrunning} 呢？
	它不是用于通信的第二个变量吗？
	确实是，而且 \apig{__sync_fetch_and_add()} 中确实隐含了
	所需的 memory barrier 指令，
	确保 thread 在检查是否应该开始之前先宣告自己的存在。
}\QuickQuizEndM
%
\QuickQuizE{
	是否有必要在访问 per-thread 变量时使用 \apik{READ_ONCE()}，
	例如，使用 \GCC 的 \apig{__thread} 存储类声明的变量？
}\QuickQuizAnswerE{
	这取决于情况。
	如果 per-thread 变量仅从其所属 thread 访问，
	且从不从信号处理程序访问，那么不需要。
	否则，很可能需要 \apik{READ_ONCE()}。
	我们将在
	\cref{sec:count:Signal-Theft Limit Counter Implementation}
	中看到这两种情况的例子。

	这引出了一个问题：一个 thread 如何能访问另一个 thread 的
	\apig{__thread} 变量？答案是第二个 thread 必须将指向其
	\apig{__thread} 变量的指针存储在第一个 thread 可以访问的
	某个位置。
	一种常见的方法是维护一个链表，每个 thread 一个元素，
	并将每个 thread 的 \apig{__thread} 变量的地址
	存储在相应的元素中。
}\QuickQuizEndE
}

\begin{fcvref}[ln:toolsoftrade:rwlockscale:reader:reader]
\clnrefrange{loop:b}{loop:e}的循环执行性能测试。
\Clnrefrange{acq:b}{acq:e}获取 lock，
\clnrefrange{hold:b}{hold:e}持有 lock 指定的微秒数，
\clnrefrange{rel:b}{rel:e}释放 lock，
\clnrefrange{think:b}{think:e}在重新获取 lock 之前等待
指定的微秒数。
\Clnref{count}计算此次 lock 获取。

\Clnref{mov_cnt}将 lock 获取计数移到该 thread 在
\co{readcounts[]} 数组中的元素，\clnref{return}返回，终止该 thread。
\end{fcvref}

\begin{figure}
\centering
\resizebox{3in}{!}{\includegraphics{CodeSamples/toolsoftrade/rwlockscale}}
\caption{在配备 Intel Xeon Platinum 8176 CPUs @ 2.10GHz 的 8 路系统上，Reader-Writer Lock 可扩展性与临界区微秒数的关系}
\label{fig:toolsoftrade:Reader-Writer Lock Scalability vs. Microseconds in Critical Section}
\end{figure}

\Cref{fig:toolsoftrade:Reader-Writer Lock Scalability vs. Microseconds in Critical Section}
展示了在一个拥有 224 核的 Xeon 系统上运行此测试的结果，
该系统每个核有两个硬件 thread，共 448 个软件可见的 CPU。
所有这些测试中 \co{thinktime} 参数为零，
\co{holdtime} 参数设置为从一微秒（图中的"1us"）到
10,000\,微秒（图中的"10000us"）的不同值。
实际绘制的值为：
\begin{equation}
	\frac{L_N}{N L_1}
\end{equation}
其中 $N$ 是当前运行中的 thread 数，
$L_N$ 是当前运行中所有 $N$ 个 thread 的 lock 获取总数，
$L_1$ 是单 thread 运行中的 lock 获取数。
在理想的硬件和软件可扩展性下，该值将始终为 1.0。

如图所示，reader-writer locking 的可扩展性明显不理想，
特别是对于较小的\IXpl{critical section}。
要理解为什么读获取可能如此缓慢，请考虑
所有获取 lock 的 thread 都必须更新 \co{pthread_rwlock_t} 数据结构。
因此，如果所有 448 个执行中的 thread 同时尝试以读模式
获取 reader-writer lock，它们必须逐个更新底层的
\co{pthread_rwlock_t}。
一个幸运的 thread 可能几乎立即完成更新，但最不幸运的
thread 必须等待其他所有 447 个 thread 完成更新。
随着 CPU 数量的增加，这种情况只会变得更糟。
还要注意 y 轴是对数刻度。
虽然 10,000\,微秒的曲线看起来相当理想，
但实际上已经比理想情况退化了约 10\pct。

\QuickQuizSeries{%
\QuickQuizB{
	与单 CPU 吞吐量比较是不是太苛刻了？
}\QuickQuizAnswerB{
	完全不是。
	事实上，如果有什么区别的话，这个比较其实过于宽松了。
	更公平的比较应该是与注释掉 locking 原语后的
	单 CPU 吞吐量进行比较。
}\QuickQuizEndB
%
\QuickQuizM{
	一微秒对于 critical section 来说并不是特别小。
	如果我需要更小的 critical section，
	例如只包含几条指令的 critical section，该怎么办？
}\QuickQuizAnswerM{
	如果被读取的数据\emph{永远}不会改变，那么你在访问它时
	不需要持有任何 lock。
	如果数据变化的频率足够低，你可以检查点执行状态，
	终止所有 thread，更改数据，然后从检查点重新启动。

	另一种方法是每个 thread 保持一个 exclusive lock，
	这样 thread 通过获取自己的 lock 来以读模式获取
	更大的聚合 reader-writer lock，通过获取所有 per-thread lock
	来以写模式获取~\cite{WilsonCHsieh92a}。
	这对 reader 来说效果很好，但会导致 writer 随着 thread 数量
	的增加承受越来越大的开销。

	\cref{chp:Deferred Processing}中描述了一些其他
	高效处理非常小的 critical section 的方法。
}\QuickQuizEndM
%
\QuickQuizE{
	所使用的系统已经有几年了，新硬件应该更快。
	那为什么还要担心 reader-writer lock 的速度呢？
}\QuickQuizAnswerE{
	总的来说，较新的硬件确实在改进。
	然而，它需要改进几个数量级才能让 reader-writer lock
	在 448 个 CPU 上达到理想性能。
	更糟糕的是，CPU 数量越多，所需的性能改进就越大。
	因此，reader-writer locking 的性能问题很可能
	在相当长的时间内伴随着我们。
}\QuickQuizEndE
}

尽管有这些限制，reader-writer locking 在许多情况下仍然非常有用，
例如当 reader 必须执行高延迟的文件或网络 I/O 时。
还有一些替代方案，其中一些将在
\cref{chp:Counting,chp:Deferred Processing}中介绍。

\subsection{原子操作（\GCC\ 经典版）}
\label{sec:toolsoftrade:Atomic Operations (gcc Classic)}

\Cref{fig:toolsoftrade:Reader-Writer Lock Scalability vs. Microseconds in Critical Section}
表明 reader-writer locking 的开销对于最小的 critical section 最为严重，
因此如果有其他方式来保护微小的 critical section 就好了。
一种这样的方式使用 \IX{atomic} 操作。
\begin{fcvref}[ln:toolsoftrade:rwlockscale:reader:reader]
我们已经见过一个 atomic 操作了，即
\cref{lst:toolsoftrade:Measuring Reader-Writer Lock Scalability}
\clnref{atmc_inc}上的 \apig{__sync_fetch_and_add()} 原语。
该原语将其第二个参数的值原子地加到其第一个参数所引用的值上，
并返回旧值（在本例中被忽略了）。
如果一对 thread 同时对同一个变量执行 \co{__sync_fetch_and_add()}，
该变量的结果值将包含两次加法的结果。
\end{fcvref}

\GNUC\ 编译器还提供了许多其他 atomic 操作，
包括 \apig{__sync_fetch_and_sub()}、
\apig{__sync_fetch_and_or()}、
\apig{__sync_fetch_and_and()}、
\apig{__sync_fetch_and_xor()} 和
\apig{__sync_fetch_and_nand()}，这些都返回旧值。
如果你需要新值，可以改用
\apig{__sync_add_and_fetch()}、
\apig{__sync_sub_and_fetch()}、
\apig{__sync_or_and_fetch()}、
\apig{__sync_and_and_fetch()}、
\apig{__sync_xor_and_fetch()} 和
\apig{__sync_nand_and_fetch()} 原语。

\QuickQuiz{
	真的有必要同时拥有两组原语吗？
}\QuickQuizAnswer{
	严格来说，不需要。
	可以用第一组的对应成员来实现第二组的任何成员。
	例如，可以用 \apig{__sync_fetch_and_nand()} 来实现
	\apig{__sync_nand_and_fetch()}，如下所示：

\begin{VerbatimU}
tmp = v;
ret = __sync_fetch_and_nand(p, tmp);
ret = ~ret & tmp;
\end{VerbatimU}

	类似地，也可以用其后值对应物来实现 \apig{__sync_fetch_and_add()}、
	\apig{__sync_fetch_and_sub()} 和 \apig{__sync_fetch_and_xor()}。

	然而，替代形式对于程序员和编译器/库实现者来说都相当方便。
}\QuickQuizEnd

经典的\IXacrml{cas}操作由一对原语提供，
即 \apig{__sync_bool_compare_and_swap()} 和
\apig{__sync_val_compare_and_swap()}。
这两个原语都原子地将一个位置更新为新值，
但仅当其先前的值等于指定的旧值时才会更新。
第一个变体在操作成功时返回 1，失败时返回 0，
例如，当先前的值不等于指定的旧值时。
第二个变体返回该位置的先前值，如果等于指定的旧值，
则表示操作成功。
\acrml{cas}操作中的任何一种都是"通用的"，
即任何对单个位置的 atomic 操作都可以用 \acrml{cas} 来实现，
尽管前面的操作在适用的地方通常更高效。
\acrml{cas}操作也能作为更广泛的 atomic 操作集的基础，
尽管其中较复杂的操作通常存在复杂性、可扩展性和性能
问题~\cite{MauriceHerlihy90a}。

\QuickQuiz{
	鉴于这些 atomic 操作通常能够生成底层指令集直接支持的
	单条 atomic 指令，它们不应该是完成任务的最快方式吗？
}\QuickQuizAnswer{
	很遗憾，不是。
	\cref{chp:Counting}中有一些鲜明的反例。
}\QuickQuizEnd

\apig{__sync_synchronize()} 原语发出一个"\IX{memory barrier}"，
它约束编译器和 CPU 重排序操作的能力，如
\cref{chp:Advanced Synchronization: Memory Ordering}中所讨论的。
在某些情况下，只需约束编译器的重排序能力而允许 CPU 自由重排序，
此时可以使用 \apik{barrier()} 原语。
\begin{fcvref}[ln:toolsoftrade:lock:reader_writer:reader]
在某些情况下，只需确保编译器不会优化掉给定的内存读取，
此时可以使用 \apik{READ_ONCE()} 原语，
如\cref{lst:toolsoftrade:Demonstration of Exclusive Locks}的
\clnref{read_x}所示。
\end{fcvref}
类似地，\apik{WRITE_ONCE()} 原语可用于防止编译器优化掉
给定的内存写入。
这最后三个原语不是由 \GCC 直接提供的，
但可以如\cref{lst:toolsoftrade:Compiler Barrier Primitive (for GCC)}
所示直接实现，
并且所有三个都在
\cref{sec:toolsoftrade:Accessing Shared Variables}中详细讨论。
另外，\apialtk{READ_ONCE(x)}{READ_ONCE()} 与 \GCC\ 内置函数
\apialtg{__atomic_load_n(&x, __ATOMIC_RELAXED)}{__atomic_load_n()}
非常相似，\apik{WRITE_ONCE()} 与 \GCC\ 内置函数
\apialtg{__atomic_store_n(&x, v, __ATOMIC_RELAXED)}{__atomic_store_n()}
非常相似。

\begin{listing}
\input{CodeSamples/api-pthreads/api-pthreads@compiler_barrier.fcv}
\caption{编译器屏障原语（适用于 \GCC）}
\label{lst:toolsoftrade:Compiler Barrier Primitive (for GCC)}
\end{listing}

\QuickQuiz{
	\apikh{ACCESS_ONCE()} 怎么了？
}\QuickQuizAnswer{
	在 2018 年的 v4.15 版本中，Linux 内核的 \apikh{ACCESS_ONCE()}
	被 \apik{READ_ONCE()} 和 \apik{WRITE_ONCE()} 分别取代，
	用于读和写操作~\cite{JonCorbet2012ACCESS:ONCE,
	JonathanCorbet2014ACCESS:ONCEcompilerBugs,
	MarkRutland2017ACCESS:ONCE:remove}。
	\apikh{ACCESS_ONCE()} 最初是作为 RCU 代码中的辅助函数引入的，
	但很快被提升为核心 API~\cite{
	PaulEMcKenney2007ACCESS:ONCE:rcu,
	LinusTorvalds2008ACCESS:ONCE:move}。
	Linux 内核的 \apik{READ_ONCE()} 和 \apik{WRITE_ONCE()} 已经
	演变为与原始 \apikh{ACCESS_ONCE()} 实现外观截然不同的复杂形式，
	这是因为需要支持大型结构的 access-once 语义，
	同时如果结构无法用单条机器指令加载和存储，
	则可能出现加载/存储撕裂。
}\QuickQuizEnd

\subsection{原子操作（C11）}
\label{sec:toolsoftrade:Atomic Operations (C11)}

C11 标准增加了 \IX{atomic} 操作，
包括加载（\apic{atomic_load()}）、
存储（\apic{atomic_store()}）、
memory barrier（\apic{atomic_thread_fence()} 和
\apic{atomic_signal_fence()}）以及读-修改-写 atomic 操作。
读-修改-写 atomic 操作包括
\apic{atomic_fetch_add()}、
\apic{atomic_fetch_sub()}、
\apic{atomic_fetch_and()}、
\apic{atomic_fetch_xor()}、
\apic{atomic_exchange()}、
\apic{atomic_compare_exchange_strong()}
和
\apic{atomic_compare_exchange_weak()}。
它们的操作方式类似于
\cref{sec:toolsoftrade:Atomic Operations (gcc Classic)}中描述的那些，
但为所有操作的 \co{_explicit} 变体增加了内存顺序参数。
在没有内存顺序参数的情况下，所有 atomic 操作都是全序的，
而参数允许更弱的顺序。
例如，``\apialtc{atomic_load_explicit(&a, memory_order_relaxed)}
{atomic_load_explicit()}''
大致类似于 Linux 内核的 ``\apik{READ_ONCE()}''。\footnote{
	内存顺序在
	\cref{chp:Advanced Synchronization: Memory Ordering}和
	\cref{chp:app:whymb:Why Memory Barriers?}中有更详细的描述。}

\subsection{原子操作（现代 \GCC）}
\label{sec:toolsoftrade:Atomic Operations (Modern gcc)}

C11 atomic 的一个限制是它们只适用于特殊的 atomic 类型，
这可能会带来问题。
因此 \GNUC\ 编译器提供了 \IX{atomic} 内置函数，包括
\apig{__atomic_load()}、
\apig{__atomic_load_n()}、
\apig{__atomic_store()}、
\apig{__atomic_store_n()}、
\apig{__atomic_thread_fence()} 等。
这些内置函数提供与其 C11 对应物相同的语义，
但可以用于普通的非 atomic 对象。
其中一些内置函数可以传递如下列表中的内存顺序参数：
\apig{__ATOMIC_RELAXED}、
\apig{__ATOMIC_CONSUME}、
\apig{__ATOMIC_ACQUIRE}、
\apig{__ATOMIC_RELEASE}、
\apig{__ATOMIC_ACQ_REL} 和
\apig{__ATOMIC_SEQ_CST}。

\subsection{Per-Thread 变量}
\label{sec:toolsoftrade:Per-Thread Variables}

Per-thread 变量，也称为线程特定数据（thread-specific data）、
线程本地存储（thread-local storage）以及其他不太礼貌的名称，
在并发代码中被大量使用，这将在
\cref{chp:Counting,chp:Data Ownership}中探讨。
POSIX 提供 \apipx{pthread_key_create()} 函数来创建
per-thread 变量（并返回相应的键），
\apipx{pthread_key_delete()} 来删除与键对应的 per-thread 变量，
\apipx{pthread_setspecific()} 来设置当前 thread 中
与指定键对应的变量值，
\apipx{pthread_getspecific()} 来返回该值。

许多编译器（包括 \GCC）提供了一个 \apig{__thread} 修饰符，
可以在变量定义中使用，将该变量指定为 per-thread 的。
然后可以正常使用变量名来访问当前 thread 的该变量实例的值。
当然，\apig{__thread} 比 POSIX 线程特定数据更容易使用，
因此对于仅使用 \GCC 或其他支持 \co{__thread} 的编译器
构建的代码，通常优先使用 \co{__thread}。

幸运的是，C11 标准引入了 \apic{_Thread_local} 关键字，
可以代替 \apig{__thread} 使用。
假以时日，这个新关键字应该能将 \co{__thread} 的易用性
与 POSIX 线程特定数据的可移植性结合起来。

\section{POSIX 操作的替代方案}
\label{sec:toolsoftrade:Alternatives to POSIX Operations}
%
\epigraph{开源的战略营销范式是一种经过达尔文主义过程过滤的
	  大规模并行的醉汉漫步。}
	 {Bruce Perens}

不幸的是，在各标准委员会着手处理 thread 操作、locking 原语和
atomic 操作之前，这些操作早已被广泛使用。
因此，这些操作的支持方式存在相当大的差异。
用汇编语言实现这些操作仍然很常见，
要么出于历史原因，要么是为了在特殊情况下获得更好的性能。
例如，\GCC 的 \co{__sync_} 系列原语都提供完整的内存顺序语义，
这在过去促使许多开发者为不需要完整内存顺序语义的情况
创建自己的实现。
以下各节展示了 Linux 内核的一些替代方案，
以及本书示例代码使用的一些历史原语。

\subsection{组织与初始化}
\label{sec:toolsoftrade:Organization and Initialization}

虽然许多环境不需要任何特殊的初始化代码，
但本书中的代码示例以调用 \apipf{smp_init()} 开始，
它初始化了从 \apipx{pthread_t} 到连续整数的映射。
用户空间 RCU 库\footnote{
	有关 RCU 的更多信息，请参阅
	\cref{sec:defer:Read-Copy Update (RCU)}。}
同样需要调用 \apiur{rcu_init()}。
虽然这些调用可以在支持构造函数的环境（如 \GCC）中隐藏，
但用户空间 RCU 库支持的大多数 RCU 风格还要求每个 thread
在创建时调用 \apiur{rcu_register_thread()}，
在退出前调用 \apiur{rcu_unregister_thread()}。

对于 Linux 内核，关于内核是否不需要调用特殊初始化代码，
还是内核的启动时代码实际上就是所需的初始化代码，
这是一个哲学问题。

\subsection{线程创建、销毁与控制}
\label{sec:toolsoftrade:Thread Creation; Destruction; and Control}

Linux 内核使用
\apik{struct task_struct} 指针来跟踪 kthread，
\apik{kthread_create()} 来创建它们，
\apik{kthread_should_stop()} 来从外部建议它们停止
（POSIX 没有对应物），\footnote{
	POSIX 环境可以通过使用适当同步的布尔标志
	配合 \co{pthread_join()} 来解决缺少
	\apik{kthread_should_stop()} 的问题。}
\apik{kthread_stop()} 来等待它们停止，
\apik{schedule_timeout_interruptible()} 用于定时等待。
还有更多的 kthread 管理 API，但这提供了一个良好的起点，
以及好的搜索关键词。

CodeSamples API 关注的是"thread"，即控制流的载体。\footnote{
	类似的软件构造还有许多其他名称，包括
	"进程"、"任务"、"纤程"、"事件"、"执行代理"
	等等。
	类似的设计原则适用于所有这些构造。}
每个这样的 thread 有一个类型为 \apipf{thread_id_t} 的标识符，
在任何给定时刻运行的两个 thread 不会有相同的标识符。
Thread 共享除 per-thread 本地状态之外的一切，\footnote{
	这算不算循环定义？}
包括程序计数器和栈。

thread API 如
\cref{lst:toolsoftrade:Thread API}所示，
各成员在以下部分描述。

\begin{listing}
\begin{VerbatimL}[numbers=none,xleftmargin=2pt]
int smp_thread_id(void)
thread_id_t create_thread(void *(*func)(void *), void *arg)
for_each_thread(t)
for_each_running_thread(t)
void *wait_thread(thread_id_t tid)
void wait_all_threads(void)
\end{VerbatimL}
\caption{Thread API 接口}
\label{lst:toolsoftrade:Thread API}
\end{listing}

\subsubsection{API 成员}

\begin{description}[style=nextline]
\item[\tco{create_thread()}]
\apipf{create_thread()} 原语创建一个新 thread，
新 thread 从 \apipf{create_thread()} 第一个参数指定的
函数 \co{func} 开始执行，并将 \apipf{create_thread()}
第二个参数指定的参数传递给它。
这个新创建的 thread 将在从 \co{func} 指定的起始函数返回时终止。
\apipf{create_thread()} 原语返回与新创建的子 thread
对应的 \apipf{thread_id_t}。

如果创建的 thread 数超过 \apipf{NR_THREADS}（包括运行程序时
隐式创建的那个），该原语将中止程序。
\apipf{NR_THREADS} 是一个编译时常量，可以修改，
不过某些系统可能对允许的 thread 数量有上限。

\item[\tco{smp_thread_id()}]
由于 \apipf{create_thread()} 返回的 \apipf{thread_id_t}
是系统相关的，\apipf{smp_thread_id()} 原语返回与发出请求的
thread 对应的 thread 索引。
该索引保证小于自程序启动以来曾经存在的最大 thread 数，
因此适用于位掩码、数组索引等用途。

\item[\tco{for_each_thread()}]
\apipf{for_each_thread()} 宏遍历所有存在的 thread，
包括如果创建\emph{将会}存在的所有 thread。
此宏对于处理
\cref{sec:toolsoftrade:Per-Thread Variables}中引入的
per-thread 变量很有用。

\item[\tco{for_each_running_thread()}]
\apipf{for_each_running_thread()} 宏仅遍历当前存在的 thread。
如果需要，与 thread 创建和删除的同步是调用者的责任。

\item[\tco{wait_thread()}]
\apipf{wait_thread()} 原语等待传入的 \co{thread_id_t}
所指定的 thread 完成。
这绝不会干扰指定 thread 的执行；
它只是等待该 thread。
注意 \apipf{wait_thread()} 返回相应 thread 的返回值。

\item[\tco{wait_all_threads()}]
\apipf{wait_all_threads()} 原语等待所有当前正在运行的 thread 完成。
如果需要，与 thread 创建和删除的同步是调用者的责任。
然而，该原语通常用于运行结束时的清理，
所以通常不需要这样的同步。

\end{description}

\subsubsection{使用示例}

\Cref{lst:toolsoftrade:Example Child Thread} (\path{threadcreate.c})
展示了一个类似 hello-world 的子 thread 示例。
如前所述，每个 thread 被分配了自己的栈，因此每个 thread
有自己的私有 \co{arg} 参数和 \co{myarg} 变量。
每个子 thread 在退出前简单地打印其参数和 \apipf{smp_thread_id()}。
注意\clnrefr{ln:intro:threadcreate:thread_test:return}上的
\co{return} 语句终止该 thread，
将 \co{NULL} 返回给对此 thread 调用 \apipf{wait_thread()} 的任何人。

\begin{listing}
\input{CodeSamples/intro/threadcreate@thread_test.fcv}
\caption{子 Thread 示例}
\label{lst:toolsoftrade:Example Child Thread}
\end{listing}

\begin{fcvref}[ln:intro:threadcreate:main]
父程序如\cref{lst:toolsoftrade:Example Parent Thread}所示。
它在\clnref{smp_init}调用 \co{smp_init()} 初始化 thread 系统，
在\clnrefrange{parse:b}{parse:e}解析参数，
在\clnref{announce}宣告其存在。
它在\clnrefrange{create:b}{create:e}创建指定数量的子 thread，
并在\clnref{wait}等待它们完成。
注意 \apipf{wait_all_threads()} 丢弃了 thread 的返回值，
因为在本例中它们都是 \co{NULL}，没什么意义。
\end{fcvref}

\begin{listing}
\input{CodeSamples/intro/threadcreate@main.fcv}
\caption{父 Thread 示例}
\label{lst:toolsoftrade:Example Parent Thread}
\end{listing}

\QuickQuiz{
	Linux 内核中与 \apipx{fork()} 和 \apipx{wait()} 等价的东西呢？
}\QuickQuizAnswer{
	它们实际上并不存在。
	在 Linux 内核中执行的所有任务都共享内存，
	除非你愿意手动做大量的内存映射工作。
}\QuickQuizEnd

\subsection{加锁}
\label{sec:toolsoftrade:Locking}

Linux 内核 locking API 的一个良好起始子集如
\cref{lst:toolsoftrade:Locking API}所示，
每个 API 元素在以下部分描述。
本书的 CodeSamples locking API 密切遵循 Linux 内核的实现。

\begin{listing}
\begin{VerbatimL}[numbers=none]
void spin_lock_init(spinlock_t *sp);
void spin_lock(spinlock_t *sp);
int spin_trylock(spinlock_t *sp);
void spin_unlock(spinlock_t *sp);
\end{VerbatimL}
\caption{锁 API 接口}
\label{lst:toolsoftrade:Locking API}
\end{listing}

\subsubsection{API 成员}

\begin{description}[style=nextline]

\item[\tco{spin_lock_init()}]
\apik{spin_lock_init()} 原语初始化指定的
\apik{spinlock_t} 变量，必须在将该变量传递给任何其他
spinlock 原语之前调用。

\item[\tco{spin_lock()}]
\apik{spin_lock()} 原语获取指定的 spinlock，
如有必要，会等待直到 spinlock 可用。
在某些环境中（如 pthreads），这种等待会涉及阻塞，
而在其他环境中（如 Linux 内核），可能涉及 CPU 忙等待的自旋循环。

关键点是在任何给定时刻只有一个 thread 可以持有一个 spinlock。

\item[\tco{spin_trylock()}]
\apik{spin_trylock()} 原语获取指定的 spinlock，
但仅在它立即可用时才获取。
如果能够获取 spinlock 则返回 \co{true}，否则返回 \co{false}。

\item[\tco{spin_unlock()}]
\apik{spin_unlock()} 原语释放指定的 spinlock，
允许其他 thread 获取它。

\end{description}

% \emph{@@@ likely need to add reader-writer locking.}

\subsubsection{使用示例}

可以使用名为 \co{mutex} 的 spinlock 来保护变量
\co{counter}，如下所示：

\begin{VerbatimU}
spin_lock(&mutex);
counter++;
spin_unlock(&mutex);
\end{VerbatimU}

\QuickQuiz{
	如果变量 \co{counter} 在没有 \co{mutex} 保护的情况下
	被递增，会出现什么问题？
}\QuickQuizAnswer{
	在具有 load-store 架构的 CPU 上，递增 \co{counter}
	可能编译成如下形式：

\begin{VerbatimU}
LOAD counter,r0
INC r0
STORE r0,counter
\end{VerbatimU}

	在这样的机器上，两个 thread 可能同时加载 \co{counter} 的值，
	各自递增它，然后各自存储结果。
	那么 \co{counter} 的新值将只比之前大一，
	尽管两个 thread 各自都递增了它。
}\QuickQuizEnd

然而，\apik{spin_lock()} 和 \apik{spin_unlock()} 原语确实有
性能影响，这将在\cref{chp:Data Structures}中看到。

\subsection{访问共享变量}
\label{sec:toolsoftrade:Accessing Shared Variables}

直到 2011 年，C 标准才为共享变量的并发读/写访问定义了语义。
然而，并发 C 代码至少在二十五年前就已经在编写了~\cite{Beck85,Inman85}。
这引出了一个问题：今天的老前辈们在遥远的 C11 之前时代是怎么做的。
对这个问题的简短回答是"他们冒着危险生活"。

\begin{listing}
\begin{fcvlabel}[ln:toolsoftrade:Living Dangerously Early 1990s Style]
\begin{VerbatimL}[commandchars=\\\{\}]
ptr = global_ptr;\lnlbl{temp}
if (ptr != NULL && ptr < high_address)
	do_low(ptr);
\end{VerbatimL}
\end{fcvlabel}
\caption{1990 年代早期的冒险风格}
\label{lst:toolsoftrade:Living Dangerously Early 1990s Style}
\end{listing}

\begin{listing}
\begin{fcvlabel}[ln:toolsoftrade:C Compilers Can Invent Loads]
\begin{VerbatimL}[commandchars=\\\{\}]
if (global_ptr != NULL &&\lnlbl{if:a}
    global_ptr < high_address)\lnlbl{if:b}
	do_low(global_ptr);\lnlbl{do_low}
\end{VerbatimL}
\end{fcvlabel}
\caption{C 编译器可以发明加载}
\label{lst:toolsoftrade:C Compilers Can Invent Loads}
\end{listing}

至少如果他们使用的是 2021 年的编译器的话，他们确实会很危险。
在（比方说）1990 年代早期，编译器做的优化较少，
部分原因是编译器开发者较少，部分原因是那个时代内存相对较小。
尽管如此，问题确实出现了，如
\cref{lst:toolsoftrade:Living Dangerously Early 1990s Style}所示，
编译器有权将其转换为
\cref{lst:toolsoftrade:C Compilers Can Invent Loads}。
如你所见，
\cref{lst:toolsoftrade:Living Dangerously Early 1990s Style}
\clnrefr{ln:toolsoftrade:Living Dangerously Early 1990s Style:temp}
上的临时变量已被优化掉，因此 \co{global_ptr} 最多会被加载三次。

\QuickQuiz{
	将\cref{lst:toolsoftrade:Living Dangerously Early 1990s Style}的
	\co{global_ptr} 加载最多三次有什么问题？
}\QuickQuizAnswer{
	假设 \co{global_ptr} 初始为非 \co{NULL}，
	但另一个 thread 将 \co{global_ptr} 设置为 \co{NULL}。
	\begin{fcvref}[ln:toolsoftrade:C Compilers Can Invent Loads]
	进一步假设转换后的代码
	(\cref{lst:toolsoftrade:C Compilers Can Invent Loads})
	的\clnref{if:a}恰好在 \co{global_ptr} 被设为 \co{NULL} 之前执行，
	\clnref{if:b}恰好在之后执行。
	那么\clnref{if:a}会认为 \co{global_ptr} 非 \co{NULL}，
	\clnref{if:b}会认为它小于 \co{high_address}，
	于是\clnref{do_low}将 \co{NULL} 指针传给 \co{do_low()}，
	而 \co{do_low()} 可能并未准备好处理这种情况。
	\end{fcvref}

	本书作者在 1990 年代早期的 DYNIX/ptx 内核内存分配器中
	正好犯了这个错误。
	追踪该 bug 不仅耗费了作者一个假期周末，
	还牵扯了他的几位同事。
	简而言之，这不是一个新问题，也不太可能自行消失。
}\QuickQuizEnd

\Cref{sec:toolsoftrade:Shared-Variable Shenanigans}
描述了普通访问导致的其他问题，
\cref{sec:toolsoftrade:A Volatile Solution,%
sec:toolsoftrade:Assembling the Rest of a Solution}
描述了一些 C11 之前的解决方案。
当然，在实际可行的情况下，应该用
\cref{sec:toolsoftrade:Atomic Operations (gcc Classic)}
或（特别是）
\cref{sec:toolsoftrade:Atomic Operations (C11)}
中描述的原语来替换直接的 C 语言内存引用。
使用这些原语来避免\IXpl{data race}，即确保如果对给定变量
有多个并发 C 语言访问，则所有这些访问都是加载操作。

\subsubsection{共享变量的诡计}
\label{sec:toolsoftrade:Shared-Variable Shenanigans}
\OriginallyPublished{sec:toolsoftrade:Shared-Variable Shenanigans}{Shared-Variable Shenanigans}{Linux Weekly News}{JadeAlglave2019WhoAfraidCompiler}
%
给定执行\IXalt{plain loads and stores}{plain access}的代码，\footnote{
	即普通的加载和存储，而不是 C11 atomic、内联汇编
	或 volatile 访问。}
编译器有权假设受影响的变量既不会被任何其他 thread 访问
也不会被修改。
这一假设允许编译器执行大量变换，
包括加载撕裂、存储撕裂、
加载融合、存储融合、代码重排序、凭空加载、
凭空存储、存储转加载变换和死代码消除，
所有这些在单 thread 代码中都工作得很好。
但并发代码可能被这些变换中的每一种破坏，
即下面描述的共享变量恶作剧。

\begin{description}[labelsep=.4em]
\item[加载撕裂（Load tearing）] 发生在编译器为单次访问使用多条加载指令时。
例如，编译器理论上可以将从 \co{global_ptr} 的加载（参见
\cref{lst:toolsoftrade:Living Dangerously Early 1990s Style}的
\clnrefr{ln:toolsoftrade:Living Dangerously Early 1990s Style:temp}）
编译为一系列单字节加载。
如果另一个 thread 同时将 \co{global_ptr} 设为 \co{NULL}，
结果可能是指针除一个字节外全部为零，
从而形成一个"野指针"。
使用这样的野指针进行存储可能破坏任意内存区域，
导致罕见且难以调试的崩溃。

更糟糕的是，在（比如说）一个具有 16 位指针的 8 位系统上，
编译器可能别无选择，只能使用一对 8 位指令来访问
给定的指针。
因为 C 标准必须支持各种系统，
标准不能在一般情况下排除加载撕裂。

\item[存储撕裂（Store tearing）] 发生在编译器为单次访问使用多条存储指令时。
例如，一个 thread 可能在同一时间将 \co{0x12345678}
存储到一个四字节整数变量中，而另一个 thread 存储了 \co{0xabcdef00}。
如果编译器对任一访问使用了 16 位存储，
结果很可能是 \co{0x1234ef00}，
这对于从该整数加载的代码来说可能是相当意外的。
这也不是纯理论问题。
例如，有些 CPU 具有较小的立即数指令字段，
在这样的 CPU 上，编译器可能将 64 位存储拆分为
两个 32 位存储，以减少在寄存器中显式构造
64 位常量的开销，即使是在 64 位 CPU 上也是如此。
有关于这种情况确实在实际中发生的历史报告
（例如~\cite{KonstantinKhlebnikov2013gccstoretearing}），
还有一份近期
报告~\cite{WillDeacon2019StoreTearingReport}。\footnote{
	请注意，即使对于正确对齐且机器字大小的访问，
	这种撕裂也可能发生，而在这种特定情况下，
	即使对于 volatile 存储也是如此。
	有些人可能会认为这种行为构成编译器 bug，
	但无论如何，它说明了从编译器编写者角度来看
	存储撕裂的感知价值。
}

当然，考虑到在 32 位系统上运行使用 64 位整数的代码的可能性，
编译器在一般情况下简直别无选择只能拆分某些存储。
但对于正确对齐的机器大小存储，\apik{WRITE_ONCE()} 将
防止存储撕裂。

\begin{listing}
\begin{fcvlabel}[ln:toolsoftrade:Preventing Load Fusing]
\begin{VerbatimLL}[commandchars=\\\{\}]
while (!need_to_stop)
	do_something_quickly();
\end{VerbatimLL}
\end{fcvlabel}
\caption{招致加载融合}
\label{lst:toolsoftrade:Inviting Load Fusing}
\end{listing}

\begin{listing}
\begin{fcvlabel}[ln:toolsoftrade:C Compilers Can Fuse Loads]
\begin{VerbatimLL}[commandchars=\\\[\]]
if (!need_to_stop)
	for (;;) {\lnlbl[loop:b]
		do_something_quickly();
		do_something_quickly();
		do_something_quickly();
		do_something_quickly();
		do_something_quickly();
		do_something_quickly();
		do_something_quickly();
		do_something_quickly();
		do_something_quickly();
		do_something_quickly();
		do_something_quickly();
		do_something_quickly();
		do_something_quickly();
		do_something_quickly();
		do_something_quickly();
		do_something_quickly();
	}\lnlbl[loop:e]
\end{VerbatimLL}
\end{fcvlabel}
\caption{C 编译器可以融合加载}
\label{lst:toolsoftrade:C Compilers Can Fuse Loads}
\end{listing}

\item[加载融合（Load fusing）] 发生在编译器使用之前从给定变量加载的结果
而不重复加载时。
这种优化不仅在单 thread 代码中完全没问题，
在多 thread 代码中通常也没问题。
不幸的是，"通常"这个词隐藏了一些真正恼人的例外。

例如，假设一个实时系统需要反复调用名为
\co{do_something_quickly()} 的函数，直到变量
\co{need_to_stop} 被设置，并且编译器可以看到
\co{do_something_quickly()} 不会存储到 \co{need_to_stop}。
一种（不安全的）编码方式如
\cref{lst:toolsoftrade:Inviting Load Fusing}所示。
编译器可能合理地将该循环展开十六次，
以减少循环末尾向后分支的每次调用开销。
更糟糕的是，因为编译器知道 \co{do_something_quickly()}
不会存储到 \co{need_to_stop}，编译器完全可以
决定只检查该变量一次，产生如
\cref{lst:toolsoftrade:C Compilers Can Fuse Loads}所示的代码。
\begin{fcvref}[ln:toolsoftrade:C Compilers Can Fuse Loads]
一旦进入，\clnrefrange{loop:b}{loop:e}上的循环将永远不会退出，
无论其他 thread 向 \co{need_to_stop} 存储非零值多少次。
\end{fcvref}
其结果最好是令人惊慌，很可能还包括严重的物理损害。

\begin{listing}
\begin{fcvlabel}[ln:toolsoftrade:C Compilers Can Fuse Non-Adjacent Loads]
\begin{VerbatimLL}[commandchars=\\\[\]]
int *gp; \lnlbl[gp]

void t0(void)
{
	WRITE_ONCE(gp, &myvar); \lnlbl[wgp]
}

void t1(void)
{
	p1 = gp; \lnlbl[p1]
	do_something(p1);
	p2 = READ_ONCE(gp); \lnlbl[p2]
	if (p2) { \lnlbl[if]
		do_something_else();
		p3 = *gp; \lnlbl[p3]
	}
}
\end{VerbatimLL}
\end{fcvlabel}
\caption{C 编译器可以融合非相邻加载}
\label{lst:toolsoftrade:C Compilers Can Fuse Non-Adjacent Loads}
\end{listing}

\begin{fcvref}[ln:toolsoftrade:C Compilers Can Fuse Non-Adjacent Loads]
编译器可以跨越惊人大范围的代码来融合加载。
例如，在
\cref{lst:toolsoftrade:C Compilers Can Fuse Non-Adjacent Loads}中，
\co{t0()} 和 \co{t1()} 并发运行，\co{do_something()} 和
\co{do_something_else()} 是内联函数。
\Clnref{gp}声明了指针 \co{gp}，C 默认将其初始化为 \co{NULL}。
在某个时刻，\co{t0()} 的\clnref{wgp}将一个非 \co{NULL}
指针存储到 \co{gp}。
同时，\co{t1()} 在\clnref{p1,p2,p3}三次从 \co{gp} 加载。
考虑到\clnref{if}发现 \co{gp} 非 \co{NULL}，人们可能
希望\clnref{p3}上的解引用永远不会引发错误。
不幸的是，编译器有权融合\clnref{p1,p3}上的读取，
这意味着如果\clnref{p1}加载了 \co{NULL} 而\clnref{p2}
加载了 \co{&myvar}，\clnref{p3}可能加载 \co{NULL}，
导致错误。\footnote{
	\ppl{Will}{Deacon}报告说这在 Linux 内核中发生过。}
注意，中间的 \apik{READ_ONCE()} 并不能阻止其他两次加载被融合，
尽管三者都是从同一个变量加载。
\end{fcvref}

\QuickQuiz{
	为什么\cref{lst:toolsoftrade:C Compilers Can Fuse Non-Adjacent Loads}
	中的 \co{do_something()} 和 \co{do_something_else()}
	是否为内联函数很重要？
}\QuickQuizAnswer{
	\begin{fcvref}[ln:toolsoftrade:C Compilers Can Fuse Non-Adjacent Loads]
	因为 \co{gp} 不是静态变量，如果 \co{do_something()}
	或 \co{do_something_else()} 是单独编译的，
	编译器必须假设这两个函数中的一个或两个可能改变
	\co{gp} 的值。
	这种可能性将迫使编译器在\clnref{p3}重新加载 \co{gp}，
	从而避免 \co{NULL} 指针解引用。
	\end{fcvref}
}\QuickQuizEnd

\item[存储融合（Store fusing）] 当编译器注意到对给定变量的连续两次存储
之间没有对该变量的加载时，就可能发生。
在这种情况下，编译器有权省略第一次存储。
这在单 thread 代码中永远不是问题，
实际上在正确编写的并发代码中通常也不是问题。
毕竟，如果两次存储快速连续执行，
其他 thread 几乎没有机会加载到第一次存储的值。

\begin{listing}
\begin{fcvlabel}[ln:toolsoftrade:C Compilers Can Fuse Stores]
\begin{VerbatimLL}[commandchars=\\\[\]]
void shut_it_down(void)
{
	status = SHUTTING_DOWN; /* BUGGY!!! */\lnlbl[store:a]
	start_shutdown();
	while (!other_task_ready) /* BUGGY!!! */\lnlbl[loop:b]
		continue;\lnlbl[loop:e]
	finish_shutdown();\lnlbl[finish]
	status = SHUT_DOWN; /* BUGGY!!! */\lnlbl[store:b]
	do_something_else();
}

void work_until_shut_down(void)
{
	while (status != SHUTTING_DOWN) /* BUGGY!!! */\lnlbl[until:loop:b]
		do_more_work();\lnlbl[until:loop:e]
	other_task_ready = 1; /* BUGGY!!! */\lnlbl[other:store]
}
\end{VerbatimLL}
\end{fcvlabel}
\caption{C 编译器可以融合存储}
\label{lst:toolsoftrade:C Compilers Can Fuse Stores}
\end{listing}

然而，也有例外，例如
\cref{lst:toolsoftrade:C Compilers Can Fuse Stores}所示。
\begin{fcvref}[ln:toolsoftrade:C Compilers Can Fuse Stores]
函数 \co{shut_it_down()} 在\clnref{store:a,store:b}
对共享变量 \co{status} 进行存储，
因此假设 \co{start_shutdown()} 和 \co{finish_shutdown()}
都不访问 \co{status}，编译器可以合理地移除\clnref{store:a}上
对 \co{status} 的存储。
不幸的是，这将意味着 \co{work_until_shut_down()} 永远不会退出
其\clnref{until:loop:b,until:loop:e}的循环，
因此永远不会设置 \co{other_task_ready}，
这反过来意味着 \co{shut_it_down()} 永远不会退出
其\clnref{loop:b,loop:e}的循环，
即使编译器选择不融合\clnref{loop:b}上对
\co{other_task_ready} 的连续加载。

\cref{lst:toolsoftrade:C Compilers Can Fuse Stores}中的代码
还有更多问题，包括代码重排序。

\item[Code reordering] 是一种常见的编译技术，用于
合并公共子表达式、减少寄存器压力，以及
提高现代超标量微处理器上众多功能单元的利用率。
它也是\cref{lst:toolsoftrade:C Compilers Can Fuse Stores}
中的代码有 bug 的另一个原因。
例如，假设\clnref{until:loop:e}上的 \co{do_more_work()} 函数
不访问 \co{other_task_ready}。
那么编译器有权将\clnref{other:store}上对 \co{other_task_ready}
的赋值移到\clnref{until:loop:b}之前，
这对于希望\clnref{until:loop:e}上最后一次调用
\co{do_more_work()} 发生在\clnref{finish}的
\co{finish_shutdown()} 调用之前的人来说，将是一个巨大的失望。
\end{fcvref}

在底层硬件可以自由重排序的情况下，
阻止编译器改变访问顺序似乎毫无意义。
然而，现代机器具有\emph{精确异常}和\emph{精确中断}，
意味着任何中断或异常都会表现为发生在指令流中的特定位置。
这意味着处理程序将看到所有先前指令的效果，
但不会看到任何后续指令的效果。
因此 \apik{READ_ONCE()} 和 \apik{WRITE_ONCE()} 可用于
控制被中断代码和中断处理程序之间的通信，
独立于底层硬件提供的顺序。\footnote{
	话虽如此，各标准委员会更希望你使用 atomic 操作
	或 \apic{sig_atomic_t} 类型的变量，
	而不是 \apik{READ_ONCE()} 和 \apik{WRITE_ONCE()}。}

\item[凭空加载（Invented load）] 已在
\cref{lst:toolsoftrade:Living Dangerously Early 1990s Style,%
lst:toolsoftrade:C Compilers Can Invent Loads}
的代码中说明过了，
编译器优化掉了一个临时变量，
从而比预期更频繁地从共享变量加载。

凭空加载可能是性能隐患。
当"热" cacheline 中的变量加载被提升到
\co{if} 语句之外时，就可能出现这些隐患。
这些提升优化并不罕见，可能导致 cache miss 显著增加，
从而严重降低性能和可扩展性。

\item[凭空存储（Invented store）] 可能在多种情况下发生。
\begin{fcvref}[ln:toolsoftrade:C Compilers Can Fuse Stores]
例如，编译器在为
\cref{lst:toolsoftrade:C Compilers Can Fuse Stores}中的
\co{work_until_shut_down()} 生成代码时可能注意到
\co{other_task_ready} 没有被 \co{do_more_work()} 访问，
且在\clnref{other:store}被存储。
如果 \co{do_more_work()} 是一个复杂的内联函数，
可能需要进行寄存器溢出，此时一个有吸引力的
临时存储位置就是 \co{other_task_ready}。
毕竟，没有对它的访问，有什么害处呢？

当然，在错误的时刻对该变量进行非零存储
将导致\clnref{loop:b}上的 \co{while} 循环过早终止，
再次允许 \co{finish_shutdown()} 与 \co{do_more_work()}
并发运行。
鉴于这个 \co{while} 的全部目的似乎是防止这种并发，
这不是一件好事。
\end{fcvref}

\begin{listing}
\begin{fcvlabel}[ln:toolsoftrade:Inviting an Invented Store]
\begin{VerbatimLL}[commandchars=\\\{\}]
if (condition)
	a = 1;
else
	do_a_bunch_of_stuff(&a);
\end{VerbatimLL}
\end{fcvlabel}
\caption{招致凭空存储}
\label{lst:toolsoftrade:Inviting an Invented Store}
\end{listing}

\begin{listing}
\begin{fcvlabel}[ln:toolsoftrade:Compiler Invents an Invited Store]
\begin{VerbatimLL}[commandchars=\\\[\]]
a = 1;\lnlbl[store:uncond]
if (!condition) {
	a = 0;\lnlbl[store:cond]
	do_a_bunch_of_stuff(&a);
}
\end{VerbatimLL}
\end{fcvlabel}
\caption{编译器发明了一个被招致的存储}
\label{lst:toolsoftrade:Compiler Invents an Invited Store}
\end{listing}

将被存储的变量用作临时变量似乎很离谱，
但标准允许这样做。
然而，读者可能有理由想要一个不那么离谱的例子，
\cref{lst:toolsoftrade:Inviting an Invented Store,%
lst:toolsoftrade:Compiler Invents an Invited Store}
提供了这样的例子。

编译器为\cref{lst:toolsoftrade:Inviting an Invented Store}
生成代码时可能知道 \co{a} 的初始值为零，
这可能强烈诱使编译器通过将该代码转换为
\cref{lst:toolsoftrade:Compiler Invents an Invited Store}
来优化掉一个分支。
\begin{fcvref}[ln:toolsoftrade:Compiler Invents an Invited Store]
这里，\clnref{store:uncond}无条件地将 \co{1} 存储到 \co{a}，
然后如果 \co{condition} 未设置，
就在\clnref{store:cond}将值重置回零。
这将 if-then-else 转换为 if-then，节省了一个分支。
\end{fcvref}

\QuickQuiz{
	哎哟！
	那编译器不是几乎可以在任何时候对普通变量发明存储吗？
}\QuickQuizAnswer{
	谢天谢地，答案是不能。
	这是因为编译器被禁止引入 data race。
	在普通存储之前发明存储的情况非常特殊：
	除非代码已经存在 data race（即使没有凭空存储），
	否则任何其他实体——无论是 CPU、thread、
	信号处理程序还是中断处理程序——都不可能看到
	被发明的存储。
	而如果代码已经存在 data race，它已经引发了
	可怕的未定义行为幽灵，允许编译器生成几乎任何它想要的代码，
	不管开发者的意愿如何。

	但如果原始存储是 volatile 的，如 \apik{WRITE_ONCE()} 中那样，
	就编译器所知，可能有与该存储相关的副作用
	可以向其他 thread 发出信号，
	允许对变量进行无 data race 的访问。
	通过发明存储，编译器可能引入 data race，
	而这是不允许的。

	此外，在
	\cref{lst:toolsoftrade:Compiler Invents an Invited Store}中，
	该变量的地址被传递给了 \co{do_a_bunch_of_stuff()}。
	如果编译器可以看到该函数的定义，
	并注意到 \co{a} 在没有任何同步操作的情况下被无条件存储，
	那么编译器可以相当确定它在这种情况下没有引入 data race。

	对于 \co{volatile} 和 atomic 变量，
	编译器被明确禁止发明写入。
}\QuickQuizEnd

最后，C11 之前的编译器可能对恰好与被写入变量相邻的
不相关变量发明写入~\cite[Section 4.2]{Boehm:2005:TCI:1064978.1065042}。
这种凭空存储的变体已被禁止编译器优化引入 data race 的
规定所取缔。

\begin{listing}
\begin{fcvlabel}[ln:toolsoftrade:Inviting a Store-to-Load Conversion]
\begin{VerbatimLL}[commandchars=\\\[\]]
r1 = p;\lnlbl[load:p]
if (unlikely(r1))\lnlbl[if]
	do_something_with(r1);\lnlbl[dsw]
barrier();\lnlbl[barrier]
p = NULL;\lnlbl[null]
\end{VerbatimLL}
\end{fcvlabel}
\caption{招致存储转加载转换}
\label{lst:toolsoftrade:Inviting a Store-to-Load Conversion}
\end{listing}

\item[存储转加载转换（Store-to-load conversion）] 当编译器注意到普通存储可能实际上
不会改变内存中的值时就可能发生。
\begin{fcvref}[ln:toolsoftrade:Inviting a Store-to-Load Conversion]
例如，考虑
\cref{lst:toolsoftrade:Inviting a Store-to-Load Conversion}。
\Clnref{load:p}获取 \co{p}，但\clnref{if}上的 \qco{if} 语句
也告诉编译器开发者认为 \co{p} 通常为零。\footnote{
	\apik{unlikely()} 函数向编译器提供此提示，
	不同的编译器提供不同的 \co{unlikely()} 实现方式。}
\clnref{barrier}上的 \apik{barrier()} 语句强制编译器忘记
\co{p} 的值，但可以想象编译器选择记住该提示——或通过
反馈驱动优化获得额外的提示。
这样做会使编译器意识到\clnref{null}通常是一个
开销很大的空操作。
\end{fcvref}

\begin{listing}
\begin{fcvlabel}[ln:toolsoftrade:Compiler Converts a Store to a Load]
\begin{VerbatimLL}[commandchars=\\\[\]]
r1 = p;\lnlbl[load:p]
if (unlikely(r1))\lnlbl[if]
	do_something_with(r1);\lnlbl[dsw]
barrier();\lnlbl[barrier]
if (p != NULL)\lnlbl[if1]
	p = NULL;\lnlbl[null]
\end{VerbatimLL}
\end{fcvlabel}
\caption{编译器将存储转换为加载}
\label{lst:toolsoftrade:Compiler Converts a Store to a Load}
\end{listing}

\begin{fcvref}[ln:toolsoftrade:Compiler Converts a Store to a Load]
这样的编译器可能因此用一个检查来保护 \co{NULL} 的存储，
如\cref{lst:toolsoftrade:Compiler Converts a Store to a Load}的
\clnref{if1,null}所示。
虽然这种转换通常是可取的，但如果实际存储对于排序是必需的，
就可能出现问题。
例如，写 memory barrier（Linux 内核的 \apik{smp_wmb()}）
会对存储排序，但不会对加载排序。
这种情况可能建议使用 \apik{smp_store_release()}
而非 \apik{smp_wmb()}。
\end{fcvref}

\item[死代码消除（Dead-code elimination）] 当编译器注意到来自加载的值
从未被使用，或一个变量被存储但从未被加载时，就可能发生。
这当然可能消除对共享变量的访问，
进而破坏 memory ordering 原语，
这可能导致你的并发代码以令人惊讶的方式运行。
迄今为止的经验表明，这些意外中令人愉快的情况相对较少。
消除只存储的变量在外部代码通过符号表定位变量的情况下
特别危险：
编译器必然不知道此类外部代码的访问，
因此可能消除外部代码依赖的变量。

\end{description}

可靠的并发代码显然需要一种方法来使编译器
保留对共享内存的重要访问的数量、顺序和类型，
这是
\cref{sec:toolsoftrade:A Volatile Solution,%
sec:toolsoftrade:Assembling the Rest of a Solution}
接下来要讨论的主题。

\subsubsection{一个 Volatile 解决方案}
\label{sec:toolsoftrade:A Volatile Solution}

虽然现在 \apic{volatile} 关键字名声不好，但在 C11 和
C++11~\cite{PeteBecker2011N3242}出现之前，
它是并行程序员工具箱中不可或缺的工具。
这引出了 \co{volatile} 到底意味着什么的问题，
即使是更新版本的标准~\cite{RichardSmith2019N4800}也没有
过于精确地回答这个问题。\footnote{
	JF Bastien 在
	C++~\cite{JFBastien2018DeprecatingVolatile}中
	详尽地记录了 \co{volatile} 关键字的历史和用例。}
该版本保证"通过 \co{volatile} \co{glvalues} 的访问
严格按照抽象机器的规则求值"，
\co{volatile} 访问是副作用，
它们是四个前进保证指标之一，
而它们的确切语义是实现定义的。
也许最清晰的指导来自这个非规范性注释：

\begin{quote}
	\apic{volatile} 是对实现的一个提示，
	要求避免对该对象进行激进的优化，
	因为该对象的值可能通过实现无法检测的方式被改变。
	此外，对于某些实现，\apic{volatile} 可能表示
	需要特殊的硬件指令来访问该对象。
	有关详细语义，请参阅 6.8.1。
	一般来说，\apic{volatile} 的语义在 C++ 中
	与在 C 中的含义相同。
\end{quote}

对于编写底层代码的人来说，这段措辞可能令人放心，
但问题是编译器编写者完全可以忽略非规范性注释。
并行程序员可以改为安慰自己：编译器编写者
不想破坏设备驱动（虽然可能要经过与设备驱动开发者的
几次"坦诚而开放的"讨论之后），
而设备驱动至少施加了以下
约束~\cite{PaulEMcKenney2016P0124R6-LKMM}：

\begin{enumerate}
\item	当该访问大小和类型的机器指令可用时，
	实现被禁止对对齐的 volatile 访问进行撕裂。\footnote{
		注意，这没有规定对于具有 128 位 CAS 但没有
		128 位加载和存储的 CPU 上的 128 位加载和存储
		应该怎么做。}
	并发代码依赖此约束来避免不必要的加载和存储撕裂。
\item	实现不得对 volatile 访问的语义做任何假设，
	对于返回值的 volatile 访问，也不得对可能返回的
	值集做任何假设。\footnote{
		这由上面提到的实现定义语义强烈暗示。}
	并发代码依赖此约束来避免在其他处理器可能正在
	并发访问相关位置时不适用的优化。
\item	对齐的机器大小的非混合大小 volatile 访问
	与前后的 volatile 汇编代码序列自然交互。
	这是必要的，因为某些设备必须使用
	volatile MMIO 访问和特殊用途汇编语言指令的
	组合来访问。
	并发代码依赖此约束，以便从 volatile 访问和
	\cref{sec:toolsoftrade:Assembling the Rest of a Solution}
	中讨论的其他手段的组合中获得所需的排序属性。
\end{enumerate}

并发代码还依赖前两个约束来避免由于 data race 引起的
未定义行为，如果对给定对象的任何访问是非 atomic 或
非 volatile 的话，假设所有访问都是对齐的且为机器大小。
混合大小访问到相同位置的语义更加复杂，
暂且搁置不谈。

那么 \apic{volatile} 相对于前面的例子表现如何？

\begin{listing}
\begin{fcvlabel}[ln:toolsoftrade:Avoiding Danger, 2018 Style]
\begin{VerbatimL}[commandchars=\\\{\}]
ptr = READ_ONCE(global_ptr);\lnlbl{temp}
if (ptr != NULL && ptr < high_address)
	do_low(ptr);
\end{VerbatimL}
\end{fcvlabel}
\caption{避免危险，2018 风格}
\label{lst:toolsoftrade:Avoiding Danger; 2018 Style}
\end{listing}

在\cref{lst:toolsoftrade:Living Dangerously Early 1990s Style}的
\clnrefr{ln:toolsoftrade:Living Dangerously Early 1990s Style:temp}
上使用 \apik{READ_ONCE()} 可以避免凭空加载，
结果如\cref{lst:toolsoftrade:Avoiding Danger; 2018 Style}所示。

\begin{listing}
\begin{fcvlabel}[ln:toolsoftrade:Preventing Load Fusing]
\begin{VerbatimL}[commandchars=\\\{\}]
while (!READ_ONCE(need_to_stop))
	do_something_quickly();
\end{VerbatimL}
\end{fcvlabel}
\caption{防止加载融合}
\label{lst:toolsoftrade:Preventing Load Fusing}
\end{listing}

如\cref{lst:toolsoftrade:Preventing Load Fusing}所示，
\apik{READ_ONCE()} 也可以防止
\cref{lst:toolsoftrade:C Compilers Can Fuse Loads}中的循环展开。

\begin{listing}
\begin{fcvlabel}[ln:toolsoftrade:Preventing Store Fusing and Invented Stores]
\begin{VerbatimL}[commandchars=\\\[\]]
void shut_it_down(void)
{
	WRITE_ONCE(status, SHUTTING_DOWN); /* BUGGY!!! */\lnlbl[store:a]
	start_shutdown();
	while (!READ_ONCE(other_task_ready)) /* BUGGY!!! */\lnlbl[loop:b]
		continue;\lnlbl[loop:e]
	finish_shutdown();\lnlbl[finish]
	WRITE_ONCE(status, SHUT_DOWN); /* BUGGY!!! */\lnlbl[store:b]
	do_something_else();
}

void work_until_shut_down(void)
{
	while (READ_ONCE(status) != SHUTTING_DOWN) /* BUGGY!!! */\lnlbl[until:loop:b]
		do_more_work();\lnlbl[until:loop:e]
	WRITE_ONCE(other_task_ready, 1); /* BUGGY!!! */\lnlbl[other:store]
}
\end{VerbatimL}
\end{fcvlabel}
\caption{防止存储融合和凭空存储}
\label{lst:toolsoftrade:Preventing Store Fusing and Invented Stores}
\end{listing}

\apik{READ_ONCE()} 和 \apik{WRITE_ONCE()} 也可以用于防止
\cref{lst:toolsoftrade:C Compilers Can Fuse Stores}中所示的
存储融合和凭空存储，
结果如\cref{lst:toolsoftrade:Preventing Store Fusing and Invented Stores}
所示。
然而，这并不能防止代码重排序，这需要
\cref{sec:toolsoftrade:Assembling the Rest of a Solution}
中教授的一些额外技巧。

\begin{listing}
\begin{fcvlabel}[ln:toolsoftrade:Disinviting an Invented Store]
\begin{VerbatimL}[commandchars=\\\{\}]
if (condition)
	WRITE_ONCE(a, 1);
else
	do_a_bunch_of_stuff();
\end{VerbatimL}
\end{fcvlabel}
\caption{取消招致凭空存储}
\label{lst:toolsoftrade:Disinviting an Invented Store}
\end{listing}

最后，\apik{WRITE_ONCE()} 可以用于防止
\cref{lst:toolsoftrade:Inviting an Invented Store}中所示的
凭空存储，
结果代码如\cref{lst:toolsoftrade:Disinviting an Invented Store}所示。

总结一下，\apic{volatile} 关键字可以在加载和存储是
机器大小且正确对齐的情况下防止加载撕裂和存储撕裂。
它还可以防止加载融合、存储融合、凭空加载
和凭空存储。
然而，虽然它确实阻止编译器相互重排序 \apic{volatile} 访问，
但它无法阻止 CPU 重排序这些访问。
此外，它无法阻止编译器或 CPU 相互重排序
非 \co{volatile} 访问，或将非 \co{volatile} 访问与
\co{volatile} 访问重排序。
防止这些类型的重排序需要下一节描述的技术。

\subsubsection{组装解决方案的其余部分}
\label{sec:toolsoftrade:Assembling the Rest of a Solution}

额外的排序传统上通过使用汇编语言来提供，
例如 \GCC\ asm 指令。
奇怪的是，这些指令实际上不需要包含汇编语言，
正如\cref{lst:toolsoftrade:Compiler Barrier Primitive (for GCC)}
所示的 \apik{barrier()} 宏所例证的。

\begin{listing}
\begin{fcvlabel}[ln:toolsoftrade:Preventing C Compilers From Fusing Loads]
\begin{VerbatimL}[commandchars=\\\[\]]
while (!need_to_stop) {
	barrier(); \lnlbl[b1]
	do_something_quickly();
	barrier(); \lnlbl[b2]
}
\end{VerbatimL}
\end{fcvlabel}
\caption{防止 C 编译器融合加载}
\label{lst:toolsoftrade:Preventing C Compilers From Fusing Loads}
\end{listing}

在 \apik{barrier()} 宏中，\co{__asm__} 引入了 asm 指令，
\co{__volatile__} 防止编译器将该 asm 优化掉，
空字符串指定不生成实际指令，
最后的 \co{"memory"} 告诉编译器这个什么都不做的 asm
可以任意改变内存。
作为回应，编译器将避免将任何内存引用移过 \apik{barrier()} 宏。
这意味着\cref{lst:toolsoftrade:C Compilers Can Fuse Loads}
中所示的破坏实时性的循环展开可以通过添加 \apik{barrier()} 调用来防止，
如\cref{lst:toolsoftrade:Preventing C Compilers From Fusing Loads}的
\clnrefr{ln:toolsoftrade:Preventing C Compilers From Fusing Loads:b1,%
ln:toolsoftrade:Preventing C Compilers From Fusing Loads:b2}
所示。
这两行代码防止编译器从任一方向将 \co{need_to_stop} 的加载
推入或推过 \co{do_something_quickly()}。

\begin{listing}
\begin{fcvlabel}[ln:toolsoftrade:Preventing Reordering]
\begin{VerbatimL}[commandchars=\\\[\]]
void shut_it_down(void)
{
	WRITE_ONCE(status, SHUTTING_DOWN);
	smp_mb(); \lnlbl[mb1]
	start_shutdown();
	while (!READ_ONCE(other_task_ready))\lnlbl[loop:b]
		continue;
	smp_mb(); \lnlbl[mb2]
	finish_shutdown();
	smp_mb(); \lnlbl[mb3]
	WRITE_ONCE(status, SHUT_DOWN);
	do_something_else();
}

void work_until_shut_down(void)
{
	while (READ_ONCE(status) != SHUTTING_DOWN) {
		smp_mb(); \lnlbl[mb4]
		do_more_work();
	}
	smp_mb(); \lnlbl[mb5]
	WRITE_ONCE(other_task_ready, 1);\lnlbl[other:store]
}
\end{VerbatimL}
\end{fcvlabel}
\caption{防止重排序}
\label{lst:toolsoftrade:Preventing Reordering}
\end{listing}

然而，这并不能阻止 CPU 对引用进行重排序。
在许多情况下，这不是问题，因为硬件只能进行一定程度的重排序。
然而，在某些情况下（如
\cref{lst:toolsoftrade:C Compilers Can Fuse Stores}）
必须约束硬件。
\Cref{lst:toolsoftrade:Preventing Store Fusing and Invented Stores}
防止了存储融合和凭空存储，
\cref{lst:toolsoftrade:Preventing Reordering}
通过添加 \apik{smp_mb()} 进一步防止了剩余的重排序，位于
\begin{fcvref}[ln:toolsoftrade:Preventing Reordering]
\clnref{mb1,mb2,mb3,mb4,mb5}。
\end{fcvref}
\apik{smp_mb()} 宏类似于
\cref{lst:toolsoftrade:Compiler Barrier Primitive (for GCC)}中所示的
\apik{barrier()}，但将空字符串替换为包含完整 memory barrier
指令的字符串，例如 x86 上的 \co{"mfence"} 或 PowerPC 上的 \co{"sync"}。

\QuickQuiz{
	但完整的 memory barrier 不是非常重量级吗？
	有没有更便宜的方法来强制
	\cref{lst:toolsoftrade:Preventing Reordering}中
	所需的排序？
}\QuickQuizAnswer{
	正如通常的情况，答案是"视情况而定"。
	然而，如果只有两个 thread 访问 \co{status}
	和 \co{other_task_ready} 变量，那么
	\cref{sec:toolsoftrade:Atomic Operations}中讨论的
	\apik{smp_store_release()} 和 \apik{smp_load_acquire()}
	函数就足够了。
}\QuickQuizEnd

排序也由一些读-修改-写 atomic 操作提供，
其中一些在\cref{sec:toolsoftrade:Atomic Operations}中介绍。
在一般情况下，内存排序可能非常微妙，如
\cref{chp:Advanced Synchronization: Memory Ordering}中所讨论的。
下一节介绍内存排序的一种替代方案，即
限制甚至完全避免 data race。

\subsubsection{避免数据竞争}
\label{sec:toolsoftrade:Avoiding Data Races}

\begin{quote}
"医生，我一想到并发访问共享变量就头疼！"

"那就别并发访问共享变量！！！"
\end{quote}

医生的建议似乎没什么用，但避免并发访问共享变量的
一个经过时间考验的方法是仅在持有特定 lock 时
访问这些变量，这将在\cref{chp:Locking}中讨论。
另一种方法是仅从给定 CPU 或 thread 访问给定的"共享"变量，
这将在\cref{chp:Data Ownership}中讨论。
可以将这两种方法结合起来，例如，给定变量可能仅由
给定 CPU 或 thread 在持有特定 lock 时修改，
并且可能从该同一 CPU 或 thread 读取，
或者从其他 CPU 或 thread 在持有该同一 lock 时读取。
在所有这些情况下，对共享变量的所有访问
都可以是普通的 C 语言访问。

以下是允许对给定变量的某些访问使用普通加载和存储，
同时要求对该同一变量的其他访问使用标记
（如 \apik{READ_ONCE()} 和 \apik{WRITE_ONCE()}）的情况列表：

\begin{enumerate}
\item	一个共享变量仅由给定的拥有者 CPU 或 thread 修改，
	但被其他 CPU 或 thread 读取。
	所有存储必须使用 \apik{WRITE_ONCE()}。
	拥有者 CPU 或 thread 可以使用普通加载。
	其他所有加载必须使用 \apik{READ_ONCE()}。
\item	一个共享变量仅在持有给定 lock 时被修改，
	但被未持有该 lock 的代码读取。
	所有存储必须使用 \apik{WRITE_ONCE()}。
	持有 lock 的 CPU 或 thread 可以使用普通加载。
	其他所有加载必须使用 \apik{READ_ONCE()}。
\item	一个共享变量仅在给定拥有者 CPU 或 thread
	持有给定 lock 时被修改，
	但被其他 CPU 或 thread 或未持有该 lock 的代码读取。
	所有存储必须使用 \apik{WRITE_ONCE()}。
	拥有者 CPU 或 thread 可以使用普通加载，
	任何持有该 lock 的 CPU 或 thread 也可以。
	其他所有加载必须使用 \apik{READ_ONCE()}。
\item	一个共享变量仅由给定 CPU 或 thread
	以及在该 CPU 或 thread 上下文中运行的
	信号或中断处理程序访问。
	处理程序可以使用普通加载和存储，
	阻止处理程序被调用的代码也可以，
	即已阻塞信号和/或中断的代码。
	所有其他代码必须使用 \apik{READ_ONCE()} 和 \apik{WRITE_ONCE()}。
\item	一个共享变量仅由给定 CPU 或 thread
	以及在该 CPU 或 thread 上下文中运行的
	信号或中断处理程序访问，
	且处理程序在返回前总是恢复其写入的所有变量的值。
	处理程序可以使用普通加载和存储，
	阻止处理程序被调用的代码也可以，
	即已阻塞信号和/或中断的代码。
	所有其他代码可以使用普通加载，但必须使用 \apik{WRITE_ONCE()}
	来防止存储撕裂、存储融合和凭空存储。
\end{enumerate}

\QuickQuiz{
	如果中断或信号处理程序本身可能被中断，会怎样？
}\QuickQuizAnswer{
	那么该中断处理程序必须遵循其他被中断代码所遵循的相同规则。
	只有那些不会被自身中断的处理程序，
	或不访问与中断处理程序共享的变量的处理程序，
	才可以安全地使用普通访问，
	而且即便如此，也只有当这些变量不会被其他 CPU 或 thread
	并发访问时才行。
}\QuickQuizEnd

在大多数其他情况下，对共享变量的加载和存储必须
分别使用 \apik{READ_ONCE()} 和 \apik{WRITE_ONCE()} 或更强的操作。
但值得重申的是，\apik{READ_ONCE()} 和 \apik{WRITE_ONCE()}
除了在编译器内部之外不提供任何排序保证。
有关此类保证的信息，请参阅上面的
\cref{sec:toolsoftrade:Assembling the Rest of a Solution}或
\cref{chp:Advanced Synchronization: Memory Ordering}。

这些 data race 避免模式的许多例子在
\cref{chp:Counting}中展示。

\subsection{原子操作}
\label{sec:toolsoftrade:Atomic Operations}

Linux 内核提供了多种 \IX{atomic} 操作，
但定义在 \apik{atomic_t} 类型上的操作提供了一个良好的起点。
普通的无撕裂读取和存储分别由
\apik{atomic_read()} 和 \apik{atomic_set()} 提供。
\IX{Acquire load} 由 \apik{smp_load_acquire()} 提供，
\IX{release store} 由 \apik{smp_store_release()} 提供。

不返回值的 fetch-and-add 操作由
\apik{atomic_add()}、\apik{atomic_sub()}、\apik{atomic_inc()} 和
\apik{atomic_dec()} 等提供。
返回是否到达零的指示的 atomic 递减由
\apik{atomic_dec_and_test()} 和 \apik{atomic_sub_and_test()} 提供。
返回新值的 atomic 加法由 \apik{atomic_add_return()} 提供。
\apik{atomic_add_unless()} 和 \apik{atomic_inc_not_zero()} 提供
条件 atomic 操作，只有当 atomic 变量的原始值与指定值不同时
才会执行（例如，这些对于管理\IXalt{reference counters}{reference count}
非常方便）。

atomic 交换操作由 \apik{atomic_xchg()} 提供，
著名的\acrmf{cas}操作由 \apik{atomic_cmpxchg()} 提供。
两者都返回旧值。
Linux 内核中有许多其他 atomic RMW 原语可用，
请参阅 Linux 内核源码树中的
\path{Documentation/atomic_t.txt} 文件。\footnote{
	截至 Linux 内核 v5.11。
}

本书的 CodeSamples API 密切遵循 Linux 内核的实现。

\subsection{Per-CPU 变量}
\label{sec:toolsoftrade:Per-CPU Variables}

Linux 内核使用 \apik{DEFINE_PER_CPU()} 来定义 per-CPU 变量，
\apik{this_cpu_ptr()} 来形成对当前 CPU 的给定 per-CPU 变量实例的引用，
\apik{per_cpu()} 来访问指定 CPU 的给定 per-CPU 变量实例，
还有许多其他专用的 per-CPU 操作。

\Cref{lst:toolsoftrade:Per-Thread-Variable API}
展示了本书的 per-thread 变量 API，
它是仿照 Linux 内核的 per-CPU 变量 API 设计的。
该 API 提供了全局变量的 per-thread 等价物。
虽然严格来说该 API 不是必需的，\footnote{
	你可以改用 \apig{__thread} 或 \apic{_Thread_local}。}
但它可以为 Linux 内核代码提供良好的用户空间类比。

\begin{listing}
\begin{VerbatimL}[numbers=none]
DEFINE_PER_THREAD(type, name)
DECLARE_PER_THREAD(type, name)
per_thread(name, thread)
__get_thread_var(name)
init_per_thread(name, v)
\end{VerbatimL}
\caption{Per-Thread 变量 API}
\label{lst:toolsoftrade:Per-Thread-Variable API}
\end{listing}

\QuickQuiz{
	在不提供 per-thread 变量 API 的系统上，
	如何解决这个问题？
}\QuickQuizAnswer{
	一种方法是创建一个以 \apipf{smp_thread_id()} 为索引的数组，
	另一种方法是使用哈希表将 \apipf{smp_thread_id()}
	映射到数组索引——这实际上就是这组 API 在
	pthread 环境中所做的。

	另一种方法是父 thread 分配一个包含每个所需
	per-thread 变量字段的结构，然后在创建 thread 时
	将其传递给子 thread。
	然而，这种方法在大型系统中可能带来巨大的软件工程成本。
	要理解这一点，想象一下如果大型系统中的所有全局变量
	都必须在单个文件中声明，
	不管它们是否是 C 静态变量！
}\QuickQuizEnd

\subsubsection{API 成员}

\begin{description}[style=nextline]
\item[\tco{DEFINE_PER_THREAD()}]

\apipf{DEFINE_PER_THREAD()} 原语定义一个 per-thread 变量。
不幸的是，无法像 Linux 内核的 \apik{DEFINE_PER_CPU()} 原语
那样提供初始化器，
但有一个 \apipf{init_per_thread()} 原语允许方便的运行时初始化。

\item[\tco{DECLARE_PER_THREAD()}]
\apipf{DECLARE_PER_THREAD()} 原语是 C 意义上的声明，而非定义。
因此，\apipf{DECLARE_PER_THREAD()} 原语可以用于访问
在其他文件中定义的 per-thread 变量。

\item[\tco{per_thread()}]
\apipf{per_thread()} 原语访问指定 thread 的变量。

\item[\tco{__get_thread_var()}]
\apipf{__get_thread_var()} 原语访问当前 thread 的变量。

\item[\tco{init_per_thread()}]
\apipf{init_per_thread()} 原语将所有 thread 的指定变量实例
设置为指定值。
Linux 内核通过普通的 C 初始化来完成此操作，
依赖于巧妙使用链接器脚本和在 CPU 上线过程中执行的代码。

\end{description}

\subsubsection{使用示例}

假设我们有一个计数器，它被非常频繁地递增，
但很少被读取。
正如\cref{sec:count:Statistical Counters}中将阐明的，
使用 per-thread 变量来实现这样的计数器是有帮助的。
这样的变量可以如下定义：

\begin{VerbatimU}
DEFINE_PER_THREAD(int, counter);
\end{VerbatimU}

计数器必须如下初始化：

\begin{VerbatimU}
init_per_thread(counter, 0);
\end{VerbatimU}

一个 thread 可以如下递增其计数器实例：

\begin{VerbatimU}
p_counter = &__get_thread_var(counter);
WRITE_ONCE(*p_counter, *p_counter + 1);
\end{VerbatimU}

计数器的值就是其所有实例的总和。
因此可以如下收集计数器值的快照：

\begin{VerbatimU}
for_each_thread(t)
  sum += READ_ONCE(per_thread(counter, t));
\end{VerbatimU}

同样，使用其他机制也可以获得类似的效果，
但 per-thread 变量兼具便利性和高性能，
这将在\cref{sec:count:Statistical Counters}中更详细地展示。

\QuickQuiz{
	如果你需要在 Linux 内核中使用 per-thread（不是 per-CPU！）
	变量，该怎么办？
}\QuickQuizAnswer{
	首先，需要 per-thread 变量的可能性比你想象的要小。
	Per-CPU 变量通常可以完成 per-thread 变量的工作。
	例如，如果你只需要做加法、按位 AND、按位 OR、交换
	或比较并交换，那么
	\co{this_cpu_add()}、
	\co{this_cpu_add_return()}、
	\co{this_cpu_and()}、
	\co{this_cpu_or()}、
	\co{this_cpu_xchg()}、
	\co{this_cpu_cmpxchg()} 和
	\co{this_cpu_cmpxchg_double()}
	操作将以低开销的方式完成任务，并且相对于上下文切换、
	中断处理程序和 softirq 处理程序是 atomic 的，
	但相对于不可屏蔽中断\emph{不是}。

	其次，在禁止抢占的代码区域内（例如，由
	\co{preempt_disable()} 和 \co{preempt_enable()} 宏包围的区域），
	当前任务保证继续在当前 CPU 上执行。
	因此，在这样的区域内，对 per-CPU 变量的任何系列访问
	相对于上下文切换是 atomic 的，
	但相对于中断处理程序、softirq 处理程序和不可屏蔽中断不是。
	但请注意，运行超过几微秒的禁止抢占代码区域
	不会受到试图构建实时系统的人的欢迎。

	第三，添加到 \co{task_struct} 结构中的字段
	充当一组 per-task 变量。
	然而，有人密切关注该结构的大小，
	这些人可能会对添加任何字段的必要性提出尖锐的问题。
	因此，如果你的字段是为仅在某些内核中构建的功能添加的，
	你绝对应该将新的 \co{task_struct} 字段
	放在适当的 \co{#ifdef} 下。

	第四，也是最后，你的 per-task 变量可能
	位于其他结构中，并由已在使用的某种同步机制保护。
	例如，如果你的代码必须持有给定 lock，
	那么对该存储的访问是否可以由该 lock 保护？
	尽管这在列表的末尾，你应该首先而不是最后考虑这种可能性！
}\QuickQuizEnd

\section{选择合适的工具}
\label{sec:toolsoftrade:The Right Tool for the Job: How to Choose?}
%
\epigraph{如果你陷入困境，换换工具；这可能解放你的思维。}
	 {Paul Arden，缩略版}

作为粗略的经验法则，使用能完成任务的最简单工具。
如果可以，就简单地顺序编程。
如果那不够，尝试使用 shell 脚本来协调并行。
如果 shell 脚本的 \co{fork()}/\co{exec()} 开销
（在 Intel Core Duo 笔记本上，一个最小 C 程序大约 480 微秒）
太大，尝试使用 C 语言的 \apipx{fork()} 和 \apipx{wait()} 原语。
如果这些原语的开销（最小子进程大约 80 微秒）仍然太大，
那么你可能需要使用 POSIX threading 原语，
选择适当的 locking 和/或 atomic 操作原语。
如果 POSIX threading 原语的开销（通常低于一微秒）
仍然太大，那么可能需要
\cref{chp:Deferred Processing}中引入的原语。
当然，实际开销不仅取决于你的硬件，
更关键的是取决于你使用原语的方式。
此外，始终记住进程间通信和消息传递可以是
共享内存多 thread 执行的良好替代方案，
特别是当你的代码充分利用了
\cref{chp:Partitioning and Synchronization Design}中
提出的设计原则时。

\QuickQuiz{
	shell 通常不是使用 \apipx{vfork()} 而不是 \apipx{fork()} 吗？
}\QuickQuizAnswer{
	很可能是这样，不过，验证留给读者作为练习。
	但同时，我希望我们可以同意 \apipx{vfork()}
	是 \apipx{fork()} 的一个变体，
	这样我们可以使用 \apipx{fork()} 作为涵盖两者的通用术语。
}\QuickQuizEnd

由于并发支持是在 C 语言首次用于构建并发系统几十年之后
才被添加到 C 标准中的，因此有多种方式并发访问共享变量。
在其他条件相同的情况下，
\cref{sec:toolsoftrade:Atomic Operations (C11)}中描述的
C11 标准操作应该是你的首选。
如果你需要对给定共享变量同时使用\IXplx{plain access}{es}和
atomic 方式访问，那么
\cref{sec:toolsoftrade:Atomic Operations (Modern gcc)}中描述的
现代 \GCC\ atomic 可能很适合你。
如果你在使用经典 \GCC\ \co{__sync} API 的旧代码库上工作，
那么你应该查看
\cref{sec:toolsoftrade:Atomic Operations (gcc Classic)}
以及相关的 \GCC\ 文档。
如果你在 Linux 内核或类似的代码库上工作，
该代码库将 \apic{volatile} 关键字与内联汇编结合使用，
或者如果你需要依赖关系来提供排序，请查看
\cref{sec:toolsoftrade:Accessing Shared Variables}
以及\cref{chp:Advanced Synchronization: Memory Ordering}
中的材料。

无论你采用哪种方法，请记住随意修改多 thread 代码
是一个极其糟糕的主意，特别是考虑到共享内存并行系统
会利用你自己的傲慢和自以为是的聪明来对付你：
你越傲慢、越觉得自己聪明，
你就会在意识到自己陷入麻烦之前给自己挖一个越深的
坑~\cite{DeadlockEmpire2016}。
因此，明智的开发者会以深深的谦虚态度来对待底层并发。
同样有必要做出正确的设计选择以及正确的个别原语选择，
这将在后续章节中详细讨论。

\QuickQuizAnswersChp{qqztoolsoftrade}
